\documentclass[10pt,a4paper,twoside]{book}

\input{Style/style}

\input{Style/matlab}

\usepackage[
	left=3cm, % inner
	right=3cm, % outer
	top=2.5cm,
	bottom=2.5cm,
	%showframe,
	]{geometry}

\usepackage{pdfpages}
\includepdfset{pages=-, offset=0mm 1cm,
   pagecommand={\thispagestyle{empty}}}

\DeclareMathOperator{\Res}{Res}
\DeclareMathOperator{\Log}{Log}
\DeclareMathOperator{\loc}{loc}

%%%%%%%% PARTE ESERCIZI

\usepackage{titlesec}
\usepackage{titletoc}
\usepackage{etoolbox}

\newcommand{\Esercizio}[1]{\subsection{Esercizio~#1}}
\newcommand{\Soluzione}{\subsection{Soluzione}}

\newcommand{\Lezione}{\qquad\hspace*{\fill}\textit}
\newcommand{\Laboratorio}{\qquad\hspace*{\fill}\textit}

\newcommand{\ParteEsercizi}{\section{Esercizi}}
\newcommand{\ParteSoluzioni}{\section{Soluzioni}}

\setcounter{secnumdepth}{2} % mi assicuro che gli esercizi siano numerati
\setcounter{tocdepth}{1} % solo 1 per non avere tutti gli esercizi nel contents

\numberwithin{figure}{chapter}
\numberwithin{equation}{chapter}

\begin{document}

\frontmatter

\pagestyle{empty}

% COPERTINA

\hypertarget{mytitlepage}{} % set the hypertarget
\bookmark[dest=mytitlepage,level=chapter]{Title Page} % add the bookmark

\vspace*{\fill}

\begin{center}

	\fg{0.5}{images/logo_polimi}

	\vspace*{1.5cm}

	{\Huge \textsc{Laboratori di CFD}}\\
	
	\vspace*{1.5cm}
	
	{\large per il corso di Fluid Labs}\\
	\vspace*{0.3cm}
	{\large A.A. 2025/2026}\\

	\vspace*{5cm}

	{\large A cura di}\\
	\vspace*{0.3cm}
	{\large Matteo Bonfadini, Manuel Alfano}\\

	\vspace*{1.5cm}

	\fg{0.2}{images/logo_blackCL}
\end{center}
\vspace*{\fill}
\clearpage

% INDICE

\cleardoublepage
\pagestyle{toc}
\hypertarget{mytoc}{} % set the hypertarget
\bookmark[dest=mytoc,level=chapter]{\contentsname} % add the bookmark
\tableofcontents
\cleardoublepage

% MAIN MATTER

\pagestyle{empty}
\mainmatter


%!TEX root = ../main.tex

\chapter{Laminar oil transport in a pipeline}
\thispagestyle{empty}

\section*{Introduzione}

\subsection*{Descrizione del caso}

\begin{figure}[H]
\centering
\includegraphics[width = 0.8\textwidth]{sketch}
\caption{Rappresentazione schematica.}
\label{fig:sketch-1}
\end{figure}%
\FloatBarrier%

Il primo test è il flusso laminare di olio in una condotta (Figura \ref{fig:sketch-1}). Se la condotta è sufficientemente lunga, ci si aspetta che il flusso sia completamente sviluppato, e questo stato deve essere riprodotto nella soluzione CFD. Tuttavia, se la condizione al contorno di ingresso è un profilo di velocità uniforme, il flusso raggiungerà lo stato completamente sviluppato solo asintoticamente. Deve quindi essere introdotto un criterio pratico per identificare la regione del tubo in cui il flusso può essere trattato come completamente sviluppato. Il flusso laminare completamente sviluppato in un tubo ammette una soluzione analitica, che è il flusso di Poiseuille cilindrico:

\begin{align}
	w (r, \theta, z) &= w(z) = - \frac{1}{4\mu} \frac{dp_e}{dz} (R^2 - r^2) \quad u(r, \theta, z) = 0 \quad v(r, \theta, z) = 0 \label{eq:poiseuille-vel} \\
	\frac{dp_e}{dz} &= \text{cost} < 0 \label{eq:pressure-grad} \\
	\tau_{rz}(r, \theta, z) &= -2\mu \frac{1}{2} \left( \frac{\partial v}{\partial z} + \frac{\partial w}{\partial r} \right) = -\mu \frac{dw}{dr} = - \frac{dp_e}{dz} \frac{r}{2} \label{eq:shear-stress}
\end{align}

dove $\mu$ è la viscosità dinamica del fluido, $R$ è il raggio interno del tubo, $u, v, w$ sono le componenti di velocità lungo le direzioni $\theta, r, z$ (Figura \ref{fig:sketch-1} e Figura \ref{fig:sketch-2}), $p_e$ è la pressione in eccesso (ovvero, la parte di pressione che non è bilanciata dalla gravità), e $\tau_{rz}$ è l'unico termine di sforzo di taglio non nullo.

\subsection*{Simulazione numerica}

Viene risolta la formulazione stazionaria delle equazioni di Navier-Stokes. Le condizioni al contorno applicate sono ingresso, uscita, parete e asse (altre condizioni al contorno potrebbero essere applicate). Il carico computazionale delle simulazioni può essere ridotto sfruttando la simmetria assiale del problema, che verrà quindi simulato come 2D attraverso un dominio avente la forma di un settore circolare, come mostrato in Figura \ref{fig:sketch-2}. All'ingresso, viene imposto un profilo di velocità $z$ uniforme pari a $W_b$ (bulk velocity), mentre le componenti di velocità $r$ e $\theta$ sono poste a zero. All'uscita, la pressione in eccesso è specificata a un dato valore, che potrebbe essere zero. Le pareti sono assunte fisse e viene impostata una condizione di aderenza (no-slip) ai bordi, ovvero la velocità del fluido è zero alle pareti. Essendo nel caso laminare, lo sforzo di taglio a parete è valutato semplicemente come la velocità assiale nelle celle vicine alla parete divisa per la sua distanza dalla parete e moltiplicata per la viscosità del fluido. Questo insieme di condizioni al contorno risulta in un campo di moto in via di sviluppo che, partendo dalla distribuzione uniforme imposta alla sezione di ingresso, raggiunge una configurazione completamente sviluppata a una certa distanza a valle di essa.

\begin{figure}[H]
\centering
\includegraphics[width = 0.8\textwidth]{bcs}
\caption{Dominio e condizioni al contorno.}
\label{fig:sketch-2}
\end{figure}%
\FloatBarrier%

\subsection*{Configurazione del problema}

\begin{itemize}
	\item Diametro del tubo $= 150~\text{mm}=0.15~\text{m}$;
	\item Bulk velocity $W_b = 0.45~\text{m/s}$;
	\item Fluido $=$ petrolio greggio pesante con densità $\rho = 910~\text{kg/m}^3$ e viscosità cinematica $\nu = 3.5 \cdot 10^{-4}~\text{m}^2/\text{s}$.
\end{itemize}

\section{Quesito 1}

Verificare che le condizioni di prova indicate corrispondano a un \textbf{flusso laminare}. Questo si ottiene calcolando il numero di Reynolds di bulk $\text{Re}_b$ e confrontandolo con il valore soglia di circa 2000.

\NotionDivider

Il flusso è laminare perché 
\begin{equation}
\label{eq:Re-b-1}
\text{Re}_b = \frac{D\cdot W_b}{\nu}=\frac{0.15 \cdot 0.45}{3.5 \cdot 10^{-4}}=192.86<2000\ .
\end{equation}

\newpage

\section{Quesito 2}
    
Impostare la lunghezza del dominio, $L$, in modo tale che si possano raggiungere condizioni di \textbf{flusso completamente sviluppato} a partire dal profilo di velocità uniforme imposto all'ingresso. 
    
Nota che l'obiettivo non è determinare con precisione la lunghezza richiesta per lo sviluppo del flusso (spesso chiamata ``lunghezza d'ingresso'' o \emph{entry length}), che è, in linea di principio, infinita. Invece, lo scopo è identificare una regione del dominio che, in pratica, possa essere considerata rappresentativa dello stato completamente sviluppato e possa quindi essere utilizzata per estrarre i risultati rilevanti della soluzione (profilo di velocità assiale, profilo degli sforzi di taglio, gradiente idraulico).

\vspace{0.3em}
    
\emph{\textbf{Suggerimento:} non sottovalutare l'impatto qualitativo! A volte un criterio qualitativo è più efficace di uno quantitativo troppo astratto.}

\vspace{0.3em}

\emph{\textbf{Suggerimento:} cerca di guardare a variabili che abbiano un chiaro significato fisico e siano più direttamente comprensibili rispetto a metriche che, sebbene significative, potrebbero sembrare oscure. Ricorda che la comprensione intuitiva del flusso, corroborata dalla conoscenza della meccanica dei fluidi, è di massima importanza nella CFD applicata.}

\vspace{0.3em}

\emph{\textbf{Suggerimento:} poiché lo stato completamente sviluppato viene raggiunto asintoticamente, qualunque sia la regione del tubo (o sezione del tubo) scelta, a rigor di logica, la soluzione CFD non corrisponderà al flusso completamente sviluppato in nessun caso. Quindi, la soluzione CFD sarà rappresentativa di un flusso completamente sviluppato al netto di una data variabilità (o, se preferisci, incertezza). L'obiettivo è avere questa variabilità entro limiti accettabili per lo scopo dello studio -- la sfida chiave è capire cosa significhi ``limiti accettabili''!}

\vspace{0.3em}

\emph{\textbf{Suggerimento:} in letteratura sono riportate molte formule per stimare la lunghezza d'ingresso del flusso laminare in un tubo. Cerca di trovarne alcune e confrontale con i tuoi risultati CFD. Fai attenzione a verificare se il confronto è coerente: per esempio, il significato di ``lunghezza d'ingresso'' nelle formule letterarie è lo stesso che stai attribuendo a questo parametro nei tuoi calcoli? Inoltre, se riesci a trovare molte formule che forniscono stime diverse, non limitarti a quelle che forniscono la corrispondenza migliore con i tuoi calcoli; chiediti invece ``perché queste formule forniscono stime diverse per gli stessi parametri? Perché alcune formule concordano con i miei risultati CFD e altre no?''}

\NotionDivider

Per iniziare la prima simulazione, scegliamo arbitrariamente come lunghezza $L=2.25~\text{m}$ e costruiamo una mesh di 100 celle in direzione assiale e 10 in direzione radiale. Osservando il profilo della velocità assiale in Figura~\ref{fig:contour1} possiamo notare che le bande di velocità si stabilizzano a circa \emph{metà} della lunghezza. 

\begin{figure}[H]
\centering
\includegraphics[width = 0.85\textwidth]{beaa1}
\caption{Colourplot velocità assiale.}
\label{fig:contour1}
\end{figure}%
\FloatBarrier%

Per capire quando lo sviluppo si arresta più precisamente studiamo l'andamento riportato in Figura~\ref{fig:m1}. Possiamo notare che da circa la cella 65 in poi, la velocità assiale lungo tre linee (linea di celle lungo l'asse di simmetria, terza fila di celle, quinta fila di celle) è costante.

\begin{figure}[H]
\centering
\includegraphics[width = 0.8\textwidth]{m1}
\caption{Andamento di DWDZ.}
\label{fig:m1}
\end{figure}%
\FloatBarrier%

A prova di ciò, plottando la velocità in direzione radiale in corrispondenza delle 65esime celle in poi, notiamo che non c'è \emph{praticamente} nessun cambiamento. Si veda Figura~\ref{fig:w1}.

\begin{figure}[H]
\centering
\includegraphics[width = 0.8\textwidth]{w1}
\caption{Andamento della velocità lungo l'asse radiale.}
\label{fig:w1}
\end{figure}%
\FloatBarrier%

Dato che siamo interessati alla \emph{entry length} (e non al numero di cella), in Figura~\ref{fig:m2} abbiamo riscalato l'ascissa del grafo di DWDZ: la \emph{entry length} $L_e$ corrisponde a circa 1.4/1.5 m.

\begin{figure}[H]
\centering
\includegraphics[width = 0.8\textwidth]{m2}
\caption{Andamento di DWDZ.}
\label{fig:m2}
\end{figure}%
\FloatBarrier%

Dal punto di vista teorico, abbiamo la formula
\begin{equation}
\label{eq:l_e=2L}
L_e=c\cdot \text{Re}\cdot D\ ,
\end{equation}

dove la costante $c$ può essere uguale a 0.05 (Langhaar, 1942, \emph{Steady flow in the transition length of a straight tube}) oppure a 0.06 (Shah \& London, 1978, \emph{Laminar Flow Forced Convection in Ducts}). Essendo il numero di Reynolds non inferiore a 100, la relazione è lineare e il parametro $c$ cambia a seconda di cosa si intenda per \emph{completamente sviluppato}: si potrebbe definire il flusso sviluppato quando la velocità al centro (oppure su un'altra linea) raggiunge il 99\% (oppure il 98\%) del valore asintotico parabolico. Come suggerito, questa è un'incertezza accettabile ai fini pratici.

Nel primo caso otteniamo $L_e=1.4464$ m, mentre nel secondo, più conservativo, otteniamo $L_e=1.7357$ m. Entrambe le formule empiriche confermano la nostra analisi. Perciò, decidiamo di scegliare una lunghezza totale di circa il doppio della \emph{entry length}: $L=2.5$ m.

\section{Quesito 3}

Verificare se la soluzione CFD è \textbf{fisicamente sensata} da un punto di vista qualitativo, cioè controllare se c'è qualcosa di strano nei risultati ispezionando i grafici a colori (colour plots) o le mappe vettoriali.

\NotionDivider

Come prima cosa, riguardando la Figura~\ref{fig:contour1} e il vector plot di Figura~\ref{fig:vect-manu}, ossserviamo che il profilo di velocità non presenta anomalie: 
\begin{itemize}
    \item non ci sono variazioni improvvise ma l'andamento è regolare e continuo, fino ad assestarsi nella regione di \emph{fully developed flow};
    \item la velocità è massima al centro del tubo cioè sull'asse di simmetria e diminuisce andando verso la parete, fino ad annullarsi per le condizioni al contorno.
\end{itemize}

Pertanto, il flusso corrisponde al caso laminare.

\begin{figure}[H]
\centering
\includegraphics[width = 0.6\textwidth]{vector-manu}
\caption{Vector plot velocità.}
\label{fig:vect-manu}
\end{figure}%
\FloatBarrier%

In aggiunta, osserviamo anche il colour plot della pressione in Figura~\ref{fig:press-manu}: diminuisce gradualmente. Questo è in accordo con l'Equazione~\eqref{eq:pressure-grad}: il fluido viene \emph{spinto} proprio grazie a questa differenza di pressione.

\begin{figure}[H]
\centering
\includegraphics[width = 0.8\textwidth]{p-manu}
\caption{Colourplot pressione.}
\label{fig:press-manu}
\end{figure}%
\FloatBarrier%

\section{Quesito 4}

Valutare la convergenza numerica della soluzione CFD. Questo richiede, da un lato, di verificare che la simulazione CFD debba essere considerata convergente rispetto all'algoritmo risolutivo e, dall'altro, che la mesh sia sufficientemente fine da mantenere l'errore di discretizzazione entro limiti accettabili.

La convergenza rispetto all'algoritmo si ottiene attraverso il monitoraggio dei residui normalizzati su tutto il campo e dei valori delle sonde (\emph{probe}), impostando un adeguato numero massimo di iterazioni e un idoneo criterio di convergenza globale.

Lo studio di sensibilità della griglia (spesso chiamato \emph{grid independence study}, ma nota che nessuna soluzione è indipendente dalla griglia in senso stretto) deve essere condotto sulla base di idonei parametri target. Nota che i parametri target sono decisi dall'utente, e tutte le quantità oggetto di indagine devono essere dimostrate indipendenti dalla griglia.

\vspace{0.3em}

\emph{\textbf{Suggerimento:} come per il raggiungimento dello stato completamente sviluppato, in linea di principio la convergenza rispetto all'algoritmo e l'indipendenza dalla griglia si ottengono asintoticamente, richiedendo un numero infinito di iterazioni (criterio di convergenza globale uguale a zero) e celle con dimensioni infinitesimali. Quindi, l'obiettivo è identificare il criterio di convergenza globale e le dimensioni delle celle in modo tale che la variabilità residua (o, se preferisci, l'incertezza) sia entro limiti accettabili. Ma, ancora una volta, la sfida chiave è capire cosa significhi limiti accettabili!}

\vspace{0.3em}

\emph{\textbf{Suggerimento:} come progetti il pattern della mesh per questo caso? Perché? Una volta definita la prima mesh, come definisci le (almeno) altre due per condurre lo studio di indipendenza della griglia?}

\NotionDivider

La prima mesh, Mesh 1, come già citato ha 100 elementi lungo l'asse e 10 elementi in direzione radiale, per un totale di 1000 elementi. Il plot dei residui di Figura~\ref{fig:res1} mostra che all'aumentare del numero di iterazioni le variabili di interesse convergono, dato che il residuo raggiunge (e oltrepassa) la soglia di $10^{-3}$. Il numero totale di iterazioni è 375.

\begin{figure}[H]
\centering
\includegraphics[width = 0.4\textwidth]{res1}
\caption{Andamento dei residui.}
\label{fig:res1}
\end{figure}%
\FloatBarrier%

Successivamente, abbiamo raffinato la nostra mesh in maniera uniforme in entrambe le direzioni: Mesh 2 ha $20\times 200 = 4000$ elementi mentre Mesh 3 ha $30\times 300 = 9000$ elementi.

All'aumentare del numero di elementi aumenta anche il numero di iterazioni necessarie per raggiungere la convergenza (375 per Mesh 1, 879 per Mesh 2, 1553 per Mesh 3), mentre \textbf{il trend dei residui si mantiene simile}, il che suggerisce l'indipendenza della soluzione dalla mesh.

Inoltre, abbiamo confrontato le variabili \texttt{W1}, \texttt{P} e \texttt{V1} tra le mesh e le Figure~\ref{fig:comp-W1},~\ref{fig:comp-P},~\ref{fig:comp-V1} confermano l'indipendenza della mesh.

\begin{figure}[H]
\centering
\includegraphics[width = 0.8\textwidth]{comp-W1}
\caption{Confronto di \texttt{W1} al centro.}
\label{fig:comp-W1}
\end{figure}%
\FloatBarrier%

\begin{figure}[H]
\centering
\includegraphics[width = 0.8\textwidth]{comp-P}
\caption{Confronto di \texttt{P1} al centro.}
\label{fig:comp-P}
\end{figure}%
\FloatBarrier%

\begin{figure}[H]
\centering
\includegraphics[width = 0.8\textwidth]{comp-V1}
\caption{Confronto di \texttt{V1} al centro.}
\label{fig:comp-V1}
\end{figure}%
\FloatBarrier%

\newpage

\section{Quesito 5}

Validare la soluzione CFD nella regione completamente sviluppata rispetto alla soluzione analitica del flusso di Poiseuille. Poiché esiste una soluzione analitica per questo caso, questa validazione è chiamata \emph{benchmarking}, e il suo scopo è semplicemente quello di auto-verificare la corretta implementazione del caso, così come la convergenza numerica. La validazione della soluzione CFD deve essere fatta in termini di profili di velocità $w(r)$ e sforzo di taglio $\tau_{rz}(r)$, così come in termini di gradiente di pressione (in eccesso) $dp_e/dz$.

Quindi, la soluzione analitica presentata a lezione, e richiamata nella sezione Descrizione del caso, deve essere manipolata in modo che tutte le quantità siano espresse in funzione della velocità media (\emph{bulk}), $W_b$, che è la quantità nota nello studio CFD (poiché è il parametro imposto alla sezione di ingresso).

\NotionDivider

Per il problema in analisi esiste una soluzione analitica con la quale si può confrontare la soluzione numerica trovata.

Poichè in pratica conosciamo solamente la velocità di bulk, cerchiamo di manipolare la soluzione del flusso di Poiseuille in modo da esprimere $\frac {d p_e}{d z}$ e $w(z)$ in termini di $W_b$. Abbiamo
\[
W_b = -\frac{1}{8\mu} \frac{d p_e}{d z} R^2\ ,
\]

da cui si ricava che il \textbf{gardiente di pressione} è
\[
\frac{d p_e}{d z} = -\frac{8\mu W_b}{R^2}\ .
\]

Utilizzando questa espressione, le equazioni originali possono essere riformulate come segue:
\begin{enumerate}
    \item \textbf{Velocità assiale}: è data da
    \[
    w(r) = 2W_b \left( 1 - \frac{r^2}{R^2} \right)\ .
    \]

    \item \textbf{Velocità radiali e tangenziali}: nel flusso puramente assiale non ci sono componenti radiali o tangenziali perciò
    \[
    u(r, \theta, z) = v(r, \theta, z) = 0\ .
    \]

    \item \textbf{Sforzo di taglio}: è dato da
    \[
    \tau_{rz}(r) = \mu\frac{dw}{dr}= \frac{4\mu W_b}{R^2} r\ .
    \]
\end{enumerate}

I confronti di Figure~\ref{fig:wr},~\ref{fig:dpedz},~\ref{fig:tau} confermano la correttezza della nostra analisi numerica anche dal punto di vista teorico.

\begin{figure}[H]
\centering
\includegraphics[width = 0.6\textwidth]{wr}
\caption{Confronto della velocità assiale.}
\label{fig:wr}
\end{figure}%
\FloatBarrier%

\begin{figure}[H]
\centering
\includegraphics[width = 0.6\textwidth]{dpedz}
\caption{Confronto del gradiente di pressione.}
\label{fig:dpedz}
\end{figure}%
\FloatBarrier%

\begin{figure}[H]
\centering
\includegraphics[width = 0.6\textwidth]{tau}
\caption{Confronto degli sforzi di taglio.}
\label{fig:tau}
\end{figure}%
\FloatBarrier%



 \cleardoublepage
%!TEX root = ../main.tex

\chapter{Turbulent water transport in a hydraulically smooth pipeline}
\thispagestyle{empty}

\section*{Introduzione}

\subsection*{Descrizione del caso}

\begin{figure}[H]
\centering
\includegraphics[width = 0.5\textwidth]{miao}
\caption{Rappresentazione schematica.}
\label{fig:sketch-12}
\end{figure}%
\FloatBarrier%

Il secondo test è il flusso turbolento di acqua in una condotta idraulicamente liscia (Figura \ref{fig:sketch-12}). A causa dell'intrinseca non stazionarietà della turbolenza, le considerazioni vengono effettuate esclusivamente in termini di media temporale. 

Se la condotta è sufficientemente lunga, ci si aspetta che il flusso sia completamente sviluppato, e questo stato deve essere riprodotto nella soluzione CFD. Tuttavia, se la condizione al contorno di ingresso è un profilo di velocità uniforme, il flusso raggiungerà lo stato completamente sviluppato solo asintoticamente. Deve quindi essere introdotto un criterio pratico per identificare la regione del tubo in cui il flusso può essere trattato come completamente sviluppato. 

Le grandezze di interesse sono le perdite per attrito lungo la condotta, così come i profili radiali della velocità assiale, degli sforzi di taglio (viscosi e turbolenti) e dei parametri di turbolenza.

A differenza del caso precedente, in quello attuale non è disponibile alcuna soluzione analitica; pertanto, non è possibile alcun benchmarking. Invece, per la validazione del setup numerico, vengono utilizzati dati sperimentali e correlazioni empiriche.


\subsection*{Simulazione numerica}

In questo laboratorio, PHOENICS viene utilizzato per simulare il flusso risolvendo le equazioni di Reynolds-Averaged Navier-Stokes (RANS) accoppiate con diversi modelli di turbolenza a viscosità turbolenta, sia nell'ambito delle \emph{wall function} che in quello di \emph{low-Reynolds}.

Inizialmente, la scelta ricadrà sul modello di turbolenza $k$-$\varepsilon$ standard con l'impiego della wall function di equilibrio di Launder e Spalding (1974).

Il carico computazionale delle simulazioni può essere ridotto sfruttando l'assialsimmetria del campo di moto medio, che sarà quindi simulato come 2D attraverso un dominio avente la forma di un settore circolare, come mostrato in Figura \ref{fig:sketch-22}.

\begin{figure}[H]
\centering
\includegraphics[width = 0.8\textwidth]{bcs}
\caption{Dominio e condizioni al contorno.}
\label{fig:sketch-22}
\end{figure}%
\FloatBarrier%

Le condizioni al contorno impiegate per simulare il secondo caso di studio sono: ingresso, uscita, pareti solide e asse.
\begin{itemize}
    \item All'ingresso, viene imposto un profilo uniforme di velocità lungo $z$ di valore $W_b$, mentre le componenti di velocità lungo $r$ e $\theta$ sono specificate come nulle. Non si assume alcuna turbolenza all'ingresso; questo si ottiene impostando il parametro di intensità di turbolenza \texttt{TI} a 0.

    \item All'uscita, viene imposta una $P^*$ esterna nulla.

    \item Alle pareti solide, a causa della condizione di no-slip, la velocità del fluido è zero. 

    Nel contesto dei Volumi Finiti con disposizione staggered-grid, solo la velocità lungo $r$ viene imposta alla parete. La condizione di no-slip per la velocità $z$ viene imposta indirettamente.

    \begin{itemize}
        \item Da un lato, il flusso avvettivo della variabile $W$ attraverso le facce delle celle vicine alla parete è impostato a zero.

        \item Dall'altro, il flusso diffusivo della variabile $W$ attraverso le facce delle celle vicine alla parete (ovvero, la forza di taglio a parete) è specificato correlando lo sforzo di taglio a parete alla velocità assiale nelle celle adiacenti alla parete tramite la wall function di equilibrio di Launder e Spalding (1974).
    \end{itemize}

    \emph{\textbf{Nota:} poiché viene adottato l'approccio wall function, la distanza adimensionale dalla parete $y^+$ dei primi nodi della griglia deve essere maggiore di 30 (e generalmente inferiore a circa 130). Tale distanza adimensionale dalla parete $y^+$ può essere calcolata solo a posteriori.}
\end{itemize}

\subsection*{Configurazione del problema}

\begin{itemize}
    \item Diametro del tubo $= 0.10~\text{m}$;
    \item Bulk-mean velocity $W_b = 0.75~\text{m/s}$;
    \item Fluido $=$ acqua a $20^\circ$C, trattata come incompressibile ($\rho = 998.23~\text{kg/m}^3$, $\nu = 10^{-6}~\text{m}^2/\text{s}$).
\end{itemize}

\subsection*{Variabili PHOENICS utili}

\begin{itemize}
    \item \texttt{YPLS}$=y^+$;
    \item \texttt{STRS}$=\tau_w/\rho$;
    \item \texttt{ENUT}$=\mu_t/\rho$;
    \item \texttt{KE}$=k$;
    \item \texttt{EP}$=\varepsilon$.
\end{itemize}

\section{Quesito 1}

Verificare che le condizioni di test indicate corrispondano a un \textbf{flusso turbolento}. Questo si ottiene calcolando il numero di Reynolds medio $\text{Re}_b$ e confrontandolo con il valore soglia di circa 2000.

\NotionDivider

Il flusso è turbolento perché
\begin{equation}
\label{eq:re2}
\text{Re}_b=\frac{D\cdot W_b}{\nu}=\frac{0.1\cdot 0.75}{10^{-6}}=75000\gg 2000\ .  
\end{equation}












 
%!TEX root = ../main.tex

\section{Quesito 2}

Impostare la lunghezza del dominio, $L$, in modo tale che si possano raggiungere condizioni di \textbf{flusso completamente sviluppato} a partire dal profilo di velocità uniforme imposto all'ingresso. Si noti che l'obiettivo non è determinare con precisione la lunghezza necessaria per lo sviluppo del flusso (spesso chiamata \emph{entry length}), che è in linea di principio infinita. Invece, lo scopo è identificare una regione del dominio che, in pratica, possa essere considerata rappresentativa dello stato completamente sviluppato e possa quindi essere utilizzata per estrarre i risultati rilevanti della soluzione (profilo di velocità assiale, profilo degli sforzi di taglio, delle variabili turbolente e del gradiente idraulico).

\vspace{0.3em}
    
\emph{\textbf{Nota:} la verifica dello stato di fluido completamente sviluppato richiede l'estrazione e il calcolo in MATLAB delle variabili elencate nella prossima sezione.}

\vspace{0.3em}
    
\emph{\textbf{Nota:} da questo punto in poi, l'interesse è solo nella regione di fully developed flow}

\vspace{0.3em}
    
\emph{\textbf{Suggerimento:} si applicano i suggerimenti del Test Case 1.}

\begin{itemize}
    \item \emph{Non sottovalutare l'impatto qualitativo! A volte un criterio qualitativo è più efficace di uno quantitativo troppo astratto.}
    \item \emph{Cerca di guardare a variabili che abbiano un chiaro significato fisico e siano più direttamente comprensibili rispetto a metriche che, sebbene significative, potrebbero sembrare oscure. Ricorda che la comprensione intuitiva del flusso, corroborata dalla conoscenza della meccanica dei fluidi, è di massima importanza nella CFD applicata.}
    \item \emph{Poiché lo stato completamente sviluppato viene raggiunto asintoticamente, qualunque sia la regione del tubo (o sezione del tubo) scelta, a rigor di logica, la soluzione CFD non corrisponderà al flusso completamente sviluppato in nessun caso. Quindi, la soluzione CFD sarà rappresentativa di un flusso completamente sviluppato al netto di una data variabilità (o, se preferisci, incertezza). L'obiettivo è avere questa variabilità entro limiti accettabili per lo scopo dello studio -- la sfida chiave è capire cosa significhi ``limiti accettabili''!}
    \item \emph{In letteratura sono riportate molte formule per stimare la lunghezza d'ingresso del flusso turbolento in un tubo. Cerca di trovarne alcune e confrontale con i tuoi risultati CFD. Inoltre, se riesci a trovare molte formule che forniscono stime diverse, non limitarti a quelle che forniscono la corrispondenza migliore con i tuoi calcoli; chiediti invece ``perché queste formule forniscono stime diverse per gli stessi parametri? Perché alcune formule concordano con i miei risultati CFD e altre no?''}
\end{itemize}

\vspace{0.3em}
    
\emph{\textbf{Suggerimento:} cerca di capire in che misura la natura del fluido (laminare e turbolento) influenzi l'entry length; potrebbe essere utile esprimere l'entry length in forma adimenisonale, normalizzandola rispetto al diametro del tubo.}

\NotionDivider

La nostra prima simulazione usa $L=2$ m. Come si evince dalla Figura~\ref{fig:r2f1}, la velocità non si è ancora stabilizza, indicatore del fatto che il fluido non è ancora completamente sviluppato.

\begin{figure}[H]
\centering
\includegraphics[width = 0.7\textwidth]{r2f1}
\caption{Velocità al centro della tubazione, $L=2$ m.}
\label{fig:r2f1}
\end{figure}%
\FloatBarrier%

Il nostro secondo tentativo prevede una lunghezza $L=10$ m. La mesh è composta da 200 elementi in direzione $z$ e 15 in direzione $r$; in particolare, gli elementi più vicini al muro ($\sim r=R$) hanno una lunghezza di 2.5 mm, mentre la restante porzione è divisa equamente in 14 parti (circa 3.4 mm).

\begin{figure}[H]
\centering
\includegraphics[width = 0.7\textwidth]{r2f2}
\caption{Velocità $W$ in diverse posizioni radiali, $L=10$ m.}
\label{fig:r2f2}
\end{figure}%
\FloatBarrier%

Analizzando qualitativamente il grafico della velocità assiale in funzione di $L/D$, Figura~\ref{fig:r2f2}, si osserva che questa grandezza raggiunge un valore costante dopo circa 45/50 diametri, circa a metà della lunghezza scelta. Ciò viene confermato anche dal grafico del gradiente della velocità, Figura~\ref{fig:r2f4}.

\begin{figure}[H]
\centering
\includegraphics[width = 0.7\textwidth]{r2f4}
\caption{Gradiente di velocità \texttt{DWDZ} in diverse posizioni, $L=10$ m.}
\label{fig:r2f4}
\end{figure}%
\FloatBarrier%

La nostra intuizione trova conferma anche nel grafico della velocità in direzione radiale (Figura~\ref{fig:r2f3}): la differenza a 5, 6 o 7 metri è minima.

\begin{figure}[H]
\centering
\includegraphics[width = 0.7\textwidth]{r2f3}
\caption{Velocità $W$ in diverse posizioni assiali, $L=10$ m.}
\label{fig:r2f3}
\end{figure}%
\FloatBarrier%

Sulla base di questa analisi qualitativa, è possibile supporre che il flusso possa essere considerato completamente sviluppato oltre una distanza $L_{\text{entry}}=50D=5$ m.

Il risultato ottenuto può essere confrontato con formule \emph{semi-empiriche} riportate in letteratura. Per esempio, F. M. White nel libro \emph{Fluid Mechanics} suggerisce una stima per la lunghezza di ingresso come segue:
\begin{equation}
\label{eq:White}
L_{\text{entry}} = 4.4 D \text{Re}^{1/6}\approx 2.9~\text{m}
\end{equation}

Guardando nuovamente il profilo della velocità in direzione radiale in Figura~\ref{fig:new3}, si può notare come a 3 m la velocità non sia del tutto ancora assestata nella prima metà del raggio, mentre lo è nella seconda metà. Ciò suggerisce che, sebbene la stima sia ottimistica, può essere ampiamente sufficiente se l'interesse dell'analisi è, per esempio, studiare il boundary layer.

Stime più restrittive considerano
\begin{equation}
\label{eq:unknown}
L_{\text{entry}} = 10 D \text{Re}^{1/6}\approx 6.4~\text{m}
\end{equation}

Esistono anche vie intermedie, ma in ogni caso esse risultano tutte comparabili a quella derivata dalla nostra analisi qualitativa.

\begin{figure}[H]
\centering
\includegraphics[width = 0.7\textwidth]{new3}
\caption{Velocità $W$ in diverse posizioni assiali.}
\label{fig:new3}
\end{figure}%
\FloatBarrier%

 
%!TEX root = ../main.tex

\section{Quesito 3}

\subsection*{Estrazione dei parametri di output rilevanti e valutazione della consistenza fisica della soluzione}

\vspace{0.5em}

I parametri di output elencati di seguito devono essere estratti dalla soluzione PHOENICS tramite post-processing in MATLAB. Questi parametri sono di interesse all'interno dello stato di flusso completamente sviluppato; tuttavia, il loro calcolo è utile anche per valutare se il flusso è ragionevolmente sviluppato o meno. Pertanto, il suggerimento è di sviluppare delle routine MATLAB per il loro calcolo all'inizio del lavoro.

I parametri sono:

\begin{itemize}
    \item La distribuzione della pressione $P^*$ lungo l'asse del tubo e il valore del coefficiente di attrito, $f$.
    
    \item Il profilo della velocità media di Reynolds lungo il raggio del tubo, che può essere espresso nelle seguenti due forme adimensionali:
    \begin{itemize}
        \item $W/W_b$ contro $r/R$, essendo $R$ il raggio del tubo.
        \item $w^+$ contro $y^+$, dove $w^+=W/w_{\tau}$ e $y^+=yw_{\tau}/\nu$, essendo $w_{\tau}=\sqrt{\tau_{w}/\rho}$ la velocità di attrito e $y$ la distanza dalla parete del tubo ($y=R-r$).
    \end{itemize}

    \emph{Presentare questi risultati sotto forma di grafico log-log.}
    
    \item Il profilo degli sforzi di taglio totali $\tau(r)$ e delle sue componenti $\tau_{\text{visc}}(r)$ e $\tau_{\text{Re}}(r)$ lungo il raggio del tubo, riferendosi alle variabili adimensionali $\tau/\tau_{w}$ ($\tau_{\text{visc}}/\tau_{w}$ e $\tau_{\text{Re}}/\tau_{w}$) contro $r/R$.
    
    \vspace{0.3em}

    \emph{\textbf{Nota:} I valori di $\tau_{\text{visc}}$ e $\tau_{\text{Re}}$ sono ottenuti numericamente dal modello costitutivo per il tensore degli sforzi e dall'ipotesi di Boussinesq per il tensore degli sforzi di Reynolds, ovvero:}
    \begin{equation}
    \label{eq:tauss}
        \tau_{\text{visc}} = \mu \frac{dW}{dr}\,, \qquad  \tau_{\text{Re}} = \mu_{t}\frac{dW}{dr}\,.
    \end{equation}
    \emph{Al contrario, lo sforzo di taglio a parete non deve essere valutato dal profilo di velocità \texttt{W1} (perché?), ma facendo riferimento alla variabile di output PHOENICS \texttt{STRS}.}
    
    \item Il profilo dell'energia cinetica turbolenta, $k$, lungo il raggio del tubo, presentato in forma adimensionale come $k/W_{b}^{2}$ contro $r/R$.
    
    \item Il profilo del tasso di dissipazione turbolenta, $\varepsilon$, lungo il raggio del tubo, presentato in forma adimensionale come $\varepsilon R/{W_{b}}^{3}$ contro $r/R$.
\end{itemize}

L'output di PHOENICS e questi parametri calcolati in MATLAB devono essere ispezionati per valutare se la soluzione CFD è \textbf{fisicamente sensata} da un punto di vista qualitativo, ovvero, controllare se c'è ``qualcosa di strano / ovviamente errato'' nei risultati.

\NotionDivider

Come prima cosa, è bene specificare che in questa simulazione $y^+$ è compreso tra 45 e 46, in accordo con il \emph{wall function} approach.

\begin{itemize}
\item Per quanto riguarda $P^*$, cioè \texttt{P1}, si osservi la Figura~\ref{fig:r2f5}. In accordo con la teoria, la pressione diminuisce linearmente.

\begin{figure}[H]
\centering
\includegraphics[width = 0.7\textwidth]{r2f5}
\caption{Pressione $P^*$ lungo l'asse $z$.}
\label{fig:r2f5}
\end{figure}%
\FloatBarrier%

La relazione teorica che lega $P^*$ e $\tau_w$ è
\begin{equation}
\label{eq:tau-P}
\boxed{
\tau_w=\frac{dP^*}{dz}\,\frac{D}{4} 
}
\end{equation}

Nel nostro caso si esprime come
\begin{equation}
\label{eq:tau-P-star}
\tau_w=\frac{P^*}{L}\,\frac{D}{4}  
\end{equation}

In PHOENICS possiamo sfruttare la variabile \texttt{STRS}$=\tau_w/\rho$ e dire che
\begin{equation}
\label{eq:tau-P-star-strs}
\texttt{STRS}= \frac{\texttt{P1}}{L}\,\frac{D}{4}\,\frac{1}{\rho} 
\end{equation}

Usando come $L=2.5$ m, corrispondente agli ultimi 2.5 m di tubo (i.e. dalla cella 150 fino alla 200), dove il fluido è completamente sviluppato, otteniamo
\begin{equation}
\label{eq:tau-P-star-calcolo}
\texttt{STRS}(\text{teorico})= \frac{134.4181}{2.5}\,\frac{0.1}{4}\,\frac{1}{998.23}=0.001347\text{ m}^2\text{s}^{-2}
\end{equation}

Invece, ispezionando la variabile \texttt{STRS} su MATLAB alla cella 150 otteniamo che \texttt{STRS}(pratico)$=0.001343\ {\text{m}^2}{\text{s}^{-2}}$ , il che verifica la consistenza (interna) della soluzione numerica.

Infine, si può calcolare il coefficiente di attrito $f$ come
\begin{equation}
\label{eq:friction-f}
\boxed{
f=\frac{4 \tau_w}{\frac{1}{2}\rho W_b^2} 
}=8\frac{ \texttt{STRS}}{ W_b^2} = 0.019106
\end{equation}

\textbf{Nota:} tale relazione deriva dalla formula di Darcy-Weisbach:
\begin{equation}
\label{eq:Darcy}
\boxed{\frac{\Delta P^*}{L} = f \, \frac{1}{D} \, \frac{1}{2}\rho W_b^2} \quad \Longrightarrow \quad f=\frac{\Delta P^*}{L}\,\frac{D}{\frac{1}{2}\rho W_b^2}=\frac{4 \tau_w}{\frac{1}{2}\rho W_b^2} \text{ dall'Equazione}~\eqref{eq:tau-P-star}
\end{equation}

\textbf{Nota:} invece che confrontare \texttt{STRS} pratico e teorico si sarebbe potuto confrontare $f$ pratico e teorico.

\item Per quanto riguarda i due contributi della velocità media di Reynolds si osservino le Figure~\ref{fig:r2f67}. Nella figura di sinistra, possiamo notare che lungo l'asse $z$ ad una lunghezza di 7.5 m dall'inlet, la velocità è \emph{praticamente} costante; tutta la variazione di velocità è confinata nella zona vicino alla parete. Questa è una caratteristica del moto turbolento, che si differenzia da quello laminare per un grafico appunto più \emph{piatto} o, in gergo, \emph{full}. Nella figura di destra osserviamo che $y^+$ parte da oltre la soglia di 30, e nella regione vicino alla parete ($\sim y^+<300$) la \emph{log law} è rispettata essendo il grafo una linea in scala logaritmica.

\begin{figure}[H]
\centering
\includegraphics[width = 0.45\textwidth]{r2f6}
\quad
\includegraphics[width = 0.45\textwidth]{r2f7}
\caption{Contributi della velocità media di Reynolds, scala logaritmica in ascissa.}
\label{fig:r2f67}
\end{figure}%
\FloatBarrier%

\item Lo sforzo di taglio e le sue componenti sono mostrati in Figura~\ref{fig:r2f8}. 

Come prima cosa, notiamo che la componente viscosa è prossima allo zero. Questo non ci sorprende, perché dal punto di vista fisico lo stress viscoso è importante solo nel sottostrato viscoso ($y^+<5$). Usando in questo caso le \emph{wall functions} con una $y^+>30$ non riusciamo a catturarne il vero comportamento.

A causa di ciò, lo sforzo totale coincide con gli sforzi di Reynolds. Quest'ultimo, al centro della tubazione cresce linearmente e vicino alla parete crolla a zero. Ciò è fisicamente corretto per gli sforzi di Reynolds poichè alla parete $\mu_t=0$, tuttavia non lo è per lo sforzo totale: esso dovrebbe corrispondere a $\tau_w$. Dovremmo avere un picco nello stress viscoso che \emph{compensa} il crollo degli sforzi di Reynolds. Come già detto, ciò non accade per via della griglia troppo grossolana.

\begin{figure}[H]
\centering
\includegraphics[width = 0.7\textwidth]{r2f8}
\caption{Profili degli sforzi.}
\label{fig:r2f8}
\end{figure}%
\FloatBarrier%

\item Il profilo dell'energia cinetica turbolenta lungo il raggio della tubazione in Figura~\ref{fig:r2f9} mostra come l'energia associata ai movimenti turbolenti vari dal centro del flusso fino alla parete della tubazione. 

Al centro, l'energia cinetica turbolenta è generalmente bassa poiché la produzione locale di turbolenza è trascurabile (il gradiente di velocità è nullo).

Verso la parete, l'energia cinetica turbolenta tende ad aumentare, raggiungendo il massimo vicino alla parete (perché i gradienti di velocità aumentano). Tuttavia, alla parete $k$ dovrebbe essere nullo per via della non-slip condition. Ciò non accade, ma è un effetto atteso dell'utilizzo delle \emph{wall functions}, che \emph{saltano} il sottostrato viscoso dove la turbolenza si smorza.

\begin{figure}[H]
\centering
\includegraphics[width = 0.7\textwidth]{r2f9}
\caption{Profilo dell'energia cinetica turbolenta.}
\label{fig:r2f9}
\end{figure}%
\FloatBarrier%

\item Il profilo del tasso di dissipazione turbolenta $\varepsilon$, Figura~\ref{fig:r2f10},  mostra come l'energia turbolenta venga dissipata lungo il raggio della tubazione. 

Al centro del flusso, la dissipazione turbolenta $\varepsilon$ è bassa, poiché la turbolenza è più debole in questa regione. 

Avvicinandosi alla parete, si osserva una crescita iperbolica di $\varepsilon$, ma i valori del tasso di dissipazione rimangono ancora molto bassi. Questo comportamento è ancora una volta dovuto alla scelta della griglia computazionale, impedendo così di cogliere correttamente la dinamica del substrato viscoso, che è determinante per la dissipazione dell'energia turbolenta vicino alla parete.

\begin{figure}[H]
\centering
\includegraphics[width = 0.7\textwidth]{r2f10}
\caption{Profilo del tasso di dissipazione turbolenta.}
\label{fig:r2f10}
\end{figure}%
\FloatBarrier%

\end{itemize}
 
%!TEX root = ../main.tex

\section{Quesito 4}

\subsection*{Valutazione della convergenza numerica della soluzione CFD}

\vspace{0.5em}

Valutare la convergenza numerica della soluzione CFD. Ciò richiede, da un lato, di verificare che la simulazione CFD sia convergente rispetto all'algoritmo di soluzione e, dall'altro, che la soluzione sia \emph{grid independent}.

La convergenza rispetto all'algoritmo di soluzione si ottiene attraverso il monitoraggio dei residui normalizzati, la stabilità dei valori puntuali e le correzioni massime/minime applicate.

Lo studio di indipendenza dalla griglia deve essere effettuato sulla base di adeguati parametri target. Si noti che i parametri target sono decisi dall'utente, e tutte le grandezze oggetto di indagine devono dimostrarsi \emph{grid-independent} (o, meglio, che l'incertezza residua dovuta a ulteriori raffinamenti della mesh rimanga entro limiti accettabili).

\NotionDivider

\begin{figure}[H]
\centering
\includegraphics[width = 0.45\textwidth]{convM1}\quad
\includegraphics[width = 0.45\textwidth]{convM2}\\
\includegraphics[width = 0.45\textwidth]{convM3}
\caption{Andamento dei residui all'aumentare della risoluzione della mesh.}
\label{fig:convcomp}
\end{figure}%
\FloatBarrier%

\begin{figure}[H]
\centering
\includegraphics[width = 0.6\textwidth]{r2f11}
\caption{Confronto della pressione $P^*$ al centro della tubazione.}
\label{fig:r2f11}
\end{figure}%
\FloatBarrier%

\begin{figure}[H]
\centering
\includegraphics[width = 0.6\textwidth]{r2f12}
\caption{Confronto di $W/W_b$ a $z=7.5$ m.}
\label{fig:r2f12}
\end{figure}%
\FloatBarrier%

\begin{figure}[H]
\centering
\includegraphics[width = 0.6\textwidth]{r2f13}
\caption{Confronto di $w^+$ a $z=7.5$ m, ascissa in scala logaritmica.}
\label{fig:r2f13}
\end{figure}%
\FloatBarrier%

\begin{figure}[H]
\centering
\includegraphics[width = 0.6\textwidth]{r2f14}
\caption{Confronto di $\tau_{\text{visc}}/\tau_w$ a $z=7.5$ m.}
\label{fig:r2f14}
\end{figure}%
\FloatBarrier%

\begin{figure}[H]
\centering
\includegraphics[width = 0.6\textwidth]{r2f15}
\caption{Confronto di $\tau_{\text{Re}}/\tau_w$ a $z=7.5$ m.}
\label{fig:r2f15}
\end{figure}%
\FloatBarrier%

\begin{figure}[H]
\centering
\includegraphics[width = 0.6\textwidth]{r2f16}
\caption{Confronto di $f$.}
\label{fig:r2f16}
\end{figure}%
\FloatBarrier%

\begin{figure}[H]
\centering
\includegraphics[width = 0.6\textwidth]{r2f17}
\caption{Confronto di $k/W_b^2$.}
\label{fig:r2f17}
\end{figure}%
\FloatBarrier%

\begin{figure}[H]
\centering
\includegraphics[width = 0.6\textwidth]{r2f18}
\caption{Confronto di $\varepsilon R/W_b^3$.}
\label{fig:r2f18}
\end{figure}%
\FloatBarrier%

Notiamo che $f$ converge verso il valore teorico di Blausius $0.316\text{Re}^{-0.25}=0.0191$.


 
%!TEX root = ../main.tex

\section{Quesito 5}

\subsection*{Validazione della soluzione CFD}

\vspace{0.5em}

Poiché non è disponibile una soluzione analitica per questo caso, i risultati CFD vengono confrontati con i dati sperimentali e alcune formule empiriche disponibili in letteratura. In particolare, i dati sperimentali provengono dalla tesi di dottorato di Lars Even Torbergsen (1998), che ha studiato il flusso pienamente sviluppato in un tubo per l'aria a numeri di Reynolds simili a quelli del caso di test attuale. Lo scopo del confronto con i dati di riferimento è duplice. Da un lato, può essere utilizzato per verificare che sia stata fatta la scelta appropriata delle impostazioni di modellazione (qui, il modello di turbolenza); in questo caso, viene usato il termine \emph{validazione}. Dall'altro, può essere utilizzato per decidere le impostazioni di modellazione più appropriate, ad esempio il miglior modello di turbolenza all'interno di un dato set di opzioni; in questo caso, viene usato il termine \emph{calibrazione}. In questa fase, si richiede di eseguire una validazione del setup CFD.

\subsubsection*{Perdite per attrito}

Innanzitutto, la validazione viene eseguita facendo riferimento alle perdite per attrito, che sono quantificate attraverso il coefficiente di attrito \(f\). Si noti che il coefficiente di attrito \(f\), il coefficiente di attrito superficiale \(C_f\), il gradiente idraulico \(J\), il gradiente di pressione \(-dP^*/dz\) e lo sforzo di taglio a parete \(\tau_w\) sono parametri intercambiabili legati alle perdite per attrito (\textbf{dimostralo!}). Pertanto, una validazione rispetto a \(f\) è praticamente equivalente a una validazione rispetto a, ad esempio, \(\tau_w\) o \(J\).

Il fattore di attrito ottenuto dalla simulazione CFD, \(f_{\text{CFD}}\), deve essere confrontato con quanto segue:
\begin{itemize}
    \item il valore ottenuto sperimentalmente da Torbergsen (1998), pari a 0.0190.
    \item il diagramma di Moody (stima grafica).
    \item la formula di Prandtl per tubi idraulicamente lisci:
    \[
    \frac{1}{\sqrt{f}} = -2 \log_{10}\left(\frac{2.51}{\text{Re}_b \sqrt{f}}\right)\,,
    \]
    \item la formula di Haaland per tubi idraulicamente lisci:
    \[
    \frac{1}{\sqrt{f}} = -1.8 \log_{10}\left(\frac{6.9}{\text{Re}_b}\right)\,.
    \]
\end{itemize}

\emph{\textbf{Nota:} La variabilità di queste stime potrebbe essere considerata come una sorta di incertezza intrinseca nella valutazione del fattore di attrito. Questo aspetto deve essere tenuto in considerazione quando si interpretano i risultati del processo di validazione.}

\NotionDivider

Il valore trovato nella terza mesh è di circa $f_{\text{CFD}}=0.01895$. Confrontiamolo con i riferimenti:
\begin{itemize}
    \item $f_{\text{Torbergsen}}=0.0190$
    \item $f_{\text{Moody}}=0.0190$
    \item $f_{\text{Prandtl}}=0.0191186$
    \item $f_{\text{Haaland}}=0.0189$
\end{itemize}

I risultati ottenuti per il fattore di attrito dalla simulazione CFD sono in buona concordanza con i valori sperimentali e le formule empiriche presenti in letteratura. Tuttavia, è essenziale considerare l'incertezza derivante sia dalle fonti numeriche che da quelle legate alla modellizzazione. Tale incertezza deve essere analizzata e quantificata per valutare la robustezza e l'affidabilità dei risultati della simulazione:
\begin{equation*}
\begin{gathered}
f=\bar{f}\pm\sigma/\sqrt{N}\,\qquad\text{dove:}\\
\bar{f}=\frac{1}{N} \sum_{i=1}^{N} f_i =\frac{f_{\text{CFD}}+f_{\text{Torbergsen}}+f_{\text{Moody}}+f_{\text{Prandtl}}+f_{\text{Haaland}}}{5}= 0.0190\,, \\
\sigma= \sqrt{\frac{1}{N-1} \sum_{i=1}^N \left(f_i - \bar{f}\right)^2}=8.12\cdot 10^{-5}\\
\Longrightarrow f=0.0190\pm 4\cdot 10^{-5}
\end{gathered}
\end{equation*}

che rientra nel nostri risultati.

\subsubsection*{Profilo della velocità assiale media di Reynolds lungo il raggio del tubo}

Successivamente, la validazione viene eseguita facendo riferimento al profilo della velocità assiale media di Reynolds. Per queste variabili, i termini di confronto sono i seguenti:
\begin{itemize}
    \item i dati ottenuti sperimentalmente da Torbergsen (1998), presentati in forma adimensionale \(W/W_b\) in funzione di \(r/R\). In particolare, vengono forniti due serie di dati, che corrispondono a due diverse operazioni di post-elaborazione sui dati sperimentali.

    \item il profilo secondo la legge di potenza di Nikuradse:
    \[
    \frac{W}{W_0} = \left(1 - \frac{r}{R}\right)^{\frac{1}{n}}\,,
    \]
    dove \(W_0\) è la velocità assiale al centro del tubo e la potenza \(n\) è uguale al fattore di attrito elevato a \( -0,5\), cioè \(n = f^{-0.5}\). Per valutare \(n\), devono essere utilizzate le stime di \(f\) dalle formule precedenti (diagramma di Moody, formula di Prandtl, formula di Haaland). Quindi, il paragone verrà fatto attraverso tre diverse curve presenti in letteratura.

    \emph{\textbf{Nota:}}
    \begin{equation*}
    \frac{W}{W_b}=\frac{(2n+1)(n+1)}{2n^2} \left( 1- \frac{r}{R}  \right)^{1/n}\,.
    \end{equation*}

    \item la funzione di parete di Launder e Spalding (1974), espressa in forma adimensionale come \(w^+\) in funzione di \(y^+\) e definita come segue:
    \[
    w^+ =
    \begin{cases} 
        y^+ & \text{se } y^+ < 11.6\,, \\
        \frac{1}{\kappa} \ln(Ey^+) & \text{se } 11.6 \leq y^+ \leq 130\,,
    \end{cases}
    \]
    dove \(\kappa = 0.41\) è la costante di Von Karman ed \(E = 8.6\) è un coefficiente di rugosità impostato per pareti idraulicamente lisce. Confronta il profilo \(w^+(y^+)\) ottenuto dalla simulazione CFD con la stima derivata dalla funzione di parete di equilibrio nel suo intervallo di validità; si consiglia di presentare i risultati in un grafico log-log. Quali considerazioni puoi fare?
\end{itemize}

\NotionDivider

Osservando i profili di velocità di Figura~\ref{fig:prof_vel_validazione}.

\begin{figure}[H]
    \centering
    \includegraphics[width=0.8\textwidth]
    {r_vel_val.jpeg}  
    \caption{Confronto dei profili di velocità.}
    \label{fig:prof_vel_validazione}
\end{figure}

Si nota una discrepanza per \( r = 0 \) (all'asse), in particolare per i grafici che seguono la legge di Nikuradse. Analizzando la relazione proposta da questa legge, si evidenziano delle differenze rispetto al modello su cui si basa la nostra simulazione CFD. La relazione iniziale è la seguente:

\[
W = W_0 \left( 1 - \frac{r}{R} \right)^{\frac{1}{n}}
\]

(Consideriamo $n=1/\sqrt{0.0190} \approx 7$)\\
1. Derivata prima rispetto a \( r \):\\
Deriviamo entrambi i membri rispetto a \( r \):
\[
\frac{dW}{dr} = W_0 \cdot \frac{d}{dr} \left( \left(1 - \frac{r}{R}\right)^{\frac{1}{n}} \right)
\]

\[
\frac{dW}{dr} = -\frac{W_0}{nR} \left(1 - \frac{r}{R}\right)^{\frac{1}{n} - 1}
\]
2. Calcolo del limite per \( r \to 0 \):\\
Sostituendo \( r = 0 \) nella derivata:
\[
\frac{dW}{dr} \bigg|_{r \to 0} = -\frac{W_0}{nR} \left(1 - \frac{0}{R}\right)^{\frac{1}{n} - 1} = -\frac{W_0}{nR}\not\to 0
\]
3. Calcolo del limite per \( r \to R \):\\
Quando \( r \to R \), la derivata tende a infinito:
\[
\frac{dW}{dr} \bigg|_{r \to R} = -\frac{W_0}{nR} \left(1 - \frac{r}{R}\right)^{\frac{1}{n} - 1} \to +\infty
\]

Si osserva una discordanza per il caso di \( r = 0 \), poiché, secondo la teoria, la derivata dovrebbe essere pari a zero in questo punto, indicando uno sforzo viscoso nullo. La simulazione CFD, invece, si comporta coerentemente con la teoria in \( r = 0 \) (dove \( \frac{dW}{dr} = 0 \)) e in \( r = R \) (dove \( \frac{dW}{dr} \to +\infty \)).

Invece, si nota un distacco tra i risultati della simulazione CFD e i dati sperimentali, nonostante questi ultimi mostrino un comportamento coerente sia all'asse che alla parete. Pertanto, sarebbe utile approfondire l'analisi del set-up dell'esperimento.

%

La Figura~\ref{fig:w_plus_validazione} mostra la coerenza tra la simulazione CFD e la teoria di Luder \& Spolding, in quanto fa riferimento a una simulazione con una griglia di dimensioni \(1 \times 20 \times 200\) senza l'uso di una power law, che determina un valore di \( y^+ \approx 45 \). Le curve in Figura~\ref{fig:w_plus_validazione} evidenziano un andamento simile tra la simulazione CFD e la teoria, con \( y^+ \) inferiore a 130. Allontanandosi dalla parete, i valori di \( y^+ \) aumentano, ma ciò che ci interessa è che le due curve siano allineate nella regione di \( y^+ \) compresa tra 45 e 130.

\begin{figure}[H]
    \centering
    \includegraphics[width=0.8\textwidth]
    {w_plus_yplus_validazione.png}  
    \caption{Confronto del profilo $w^+ \ \text{vs} \  y^+$.}
    \label{fig:w_plus_validazione}
\end{figure}

\subsubsection*{Profilo dell'energia cinetica turbolenta lungo il raggio del tubo}

La validazione viene estesa al profilo dell'energia cinetica turbolenta lungo il raggio del tubo. Si raccomanda di presentare i risultati in forma adimensionale come \(k/W_b^2\) in funzione di \(r/R\). I termini di confronto sono i seguenti:
\begin{itemize}
    \item i dati ottenuti sperimentalmente da Torbergsen (1998), già presentati in forma adimensionale \(k/W_b^2\) in funzione di \(r/R\).
    \item la seguente formula empirica riportata nella tesi di Kam Hong Ng (1971), già presentata in forma adimensionale:
    \[
    \frac{k}{W_b^2} = \frac{f}{8} \left[1 + \frac{2}{3} \left(\frac{r}{R}\right) + \frac{10}{3} \left(\frac{r}{R}\right)^3\right]\,.
    \]
\end{itemize}
Anche in questo caso, devono essere utilizzate le stime di \(f\) dalle formule precedenti.

\NotionDivider

Osserviamo la Figura~\ref{fig:K_validazione}:
\begin{figure}[H]
    \centering
    \includegraphics[width=0.7\textwidth]
    {K_validazione.png}  
    \caption{Confronto profilo energia cinetica turbolenta.}
    \label{fig:K_validazione}
\end{figure}

I grafici ottenuti mostrano un andamento qualitativo coerente con le aspettative, con un'energia cinetica turbolenta inizialmente bassa all'asse che aumenta progressivamente fino a raggiungere un picco in prossimità della parete. Questo comportamento è attribuibile alla dominanza della turbolenza e del moto caotico a causa dell'interazione con la parete. Tuttavia, è stata osservata una discrepanza significativa tra i dati sperimentali, la simulazione CFD e la formula empirica di Kam Hong. La differenza nei dati sperimentali può essere spiegata dal fatto che la misurazione diretta dell'energia cinetica turbolenta in tale regione è particolarmente complessa, in quanto si tratta di flussi turbolenti. È plausibile che i dati sperimentali siano ottenuti indirettamente, utilizzando misurazioni di grandezze correlate.

Per comprendere meglio la discrepanza con la formula empirica di Kam Hong, sarebbe opportuno analizzare nel dettaglio la metodologia di derivazione di tale relazione.

In generale, ci si può aspettare un certo discostamento dai valori teorici man mano che ci si avvicina alla parete, in quanto la dinamica del flusso in questa regione diventa più complessa. Da un lato, è difficile per i modelli catturare correttamente questo comportamento, mentre, dall'altro lato, la misurazione sperimentale di queste grandezze è altrettanto complessa, rendendo difficile ottenere dati precisi in prossimità della parete.

\subsubsection*{Profilo del tasso di dissipazione turbolenta lungo il raggio del tubo}

Infine, la validazione si concentra sul profilo del tasso di dissipazione turbolenta lungo il raggio del tubo. Si raccomanda di presentare i risultati in forma adimensionale come \(\epsilon R/W_b^3\) in funzione di \(r/R\). I termini di confronto sono i seguenti:
\begin{itemize}
    \item i dati ottenuti sperimentalmente da Torbergsen (1998), già presentati in forma adimensionale \(\epsilon R/W_b^3\) in funzione di \(r/R\). In particolare, vengono fornite tre serie di dati, che corrispondono a metodi diversi per stimare \(\epsilon\) dalla stessa serie di dati sperimentali. La differenza tra le tre serie sottolinea le difficoltà nel fornire stime sperimentali di \(\epsilon\),
    \item la seguente relazione alla base del modello di turbolenza \(k-\epsilon\):
    \[
    \epsilon(r) = C_d \frac{k^2(r)}{l_m(r)} = 0.1643 \frac{k^2(r)}{l_m(r)},
    \]
    dove \(l_m\) è la lunghezza di miscelazione, che è una caratteristica tipica della scala di turbolenza.
\end{itemize}

\NotionDivider

\begin{figure}[H]
    \centering
    \includegraphics[width=0.7\textwidth]
    {eps_validation.png}  
    \caption{Confronto del profilo del tasso di dissipazione turbolenta.}
    \label{fig:eps_validazione}
\end{figure}

La Figura~\ref{fig:eps_validazione} mostra che i tassi di dissipazione turbolenta assumono valori relativamente bassi, a causa della condizione imposta nella simulazione CFD, in cui si considera $y^+ > 5$, escludendo lo sviluppo del sublayer viscoso. Tuttavia, in prossimità della parete, il modello turbolento evidenzia un'esplosione del tasso di dissipazione, probabilmente dovuta alla capacità del modello di catturare la dinamica del sublayer viscoso, responsabile della forte dissipazione turbolenta vicino alla parete.

Inoltre, la stessa figura evidenzia come i tassi di dissipazione turbolenta si discostino ulteriormente avvicinandosi alla parete, dove la dinamica caotica diventa sufficientemente complessa da rendere difficile la sua descrizione sia attraverso i modelli turbolenti, sia mediante misurazioni sperimentali. È probabile che queste ultime siano ottenute in modo indiretto, utilizzando misure di grandezze correlate a quella di interesse.

\newpage
 
%!TEX root = ../main.tex

\section{Quesito 6}

\subsection*{Analisi dell'influenza del modello di turbolenza}

\vspace{0.5em}

I risultati CFD potrebbero essere sensibili alla scelta del modello di turbolenza. Lo scopo di questa parte dello studio è capire quali parametri di output sono maggiormente influenzati e in che misura. Tale ``incertezza di modellazione'' deve essere presa in considerazione quando si confrontano i risultati CFD con i dati sperimentali o quando si producono stime CFD in assenza di dati sperimentali per la calibrazione.

\begin{enumerate}
    \item Inizialmente, il test verrà effettuato considerando altre due varianti popolari del modello di turbolenza $k$-$\varepsilon$ standard, tutte operanti all'interno dell'approccio \emph{wall function} (richiedendo quindi che $y^+$ dei primi nodi della griglia sia maggiore di 30 e inferiore a circa 130).
    Questi modelli sono:
    \begin{itemize}
        \item Il modello di turbolenza RNG $k$-$\varepsilon$ di Yakhot e Orszag (1986).
        \item Il modello di turbolenza Realisable $k$-$\varepsilon$ di Shih et al. (1995).
    \end{itemize}

    \item In alternativa all'approccio wall function, quello \textbf{low-Reynolds} si basa sulla risoluzione del campo di moto vicino alla parete fino al sottostrato viscoso. Ciò richiede, da un lato, la definizione di mesh molto fini vicino alla parete, in modo che le celle near-wall siano caratterizzate da $y^+ \approx 1$ e, dall'altro, la modifica delle equazioni del modello di turbolenza vicino alla parete mediante funzioni di smorzamento.
    
    Risolvere lo stesso caso di studio utilizzando un modello di turbolenza low-Reynolds, nello specifico il modello Two-Layer $k$-$\varepsilon$ di Rodi (1991).
    
    Prestare attenzione alla definizione della mesh e fare le proprie considerazioni su punti di forza e debolezze dell'approccio \emph{low-Reynolds} rispetto a quello \emph{wall function}. Ripetere l'analisi e la validazione anche per questo nuovo modello di turbolenza e, in particolare, analizzare le informazioni che si possono raccogliere sul profilo dello sforzo di taglio e le sue componenti (viscoso e turbolento) utilizzando l'approccio \emph{low Reynolds}.
\end{enumerate}

\NotionDivider

\begin{enumerate}
\item $k$-$\varepsilon$ VS RNG $k$-$\varepsilon$ VS Realisable $k$-$\varepsilon$

\begin{figure}[H]
\centering
\includegraphics[width = 0.6\textwidth]{r3f1}
\caption{Confronto della pressione $P^*$ al centro della tubazione.}
\label{fig:r3f1}
\end{figure}%
\FloatBarrier%

\begin{figure}[H]
\centering
\includegraphics[width = 0.5\textwidth]{r3f7}
\caption{Confronto di $W/W_b$ a $z=7.5$ m.}
\label{fig:r3f7}
\end{figure}%
\FloatBarrier%

\begin{figure}[H]
\centering
\includegraphics[width = 0.5\textwidth]{r3f9}
\caption{Confronto di \texttt{DWDZ}.}
\label{fig:r3f9}
\end{figure}%
\FloatBarrier%

Notiamo che la condizione di FD è raggiunta dopo per Realisable $k$-$\varepsilon$.

\begin{figure}[H]
\centering
\includegraphics[width = 0.5\textwidth]{r3f8}
\caption{Confronto di $w^+$ a $z=7.5$ m, ascissa in scala logaritmica.}
\label{fig:r3f8}
\end{figure}%
\FloatBarrier%

Qui possiamo notare che al centro del tubo (per $y^+$ alti) i modelli RNG e Realisable riescono a moellizzare una velocità più alta perché il mescolamento turbolento ($C_\mu$) non è costante ma dipende dalla deformazione del fluido (strain rate).

\newpage

\begin{figure}[H]
\centering
\includegraphics[width = 0.6\textwidth]{r3f4}
\caption{Confronto di $\tau_{\text{visc}}/\tau_w$ a $z=7.5$ m.}
\label{fig:r3f4}
\end{figure}%
\FloatBarrier%

\begin{figure}[H]
\centering
\includegraphics[width = 0.6\textwidth]{r3f5}
\caption{Confronto di $\tau_{\text{Re}}/\tau_w$ a $z=7.5$ m.}
\label{fig:r3f5}
\end{figure}%
\FloatBarrier%

\begin{figure}[H]
\centering
\includegraphics[width = 0.6\textwidth]{r3f6}
\caption{Confronto di $f$.}
\label{fig:r3f6}
\end{figure}%
\FloatBarrier%

\begin{figure}[H]
\centering
\includegraphics[width = 0.6\textwidth]{r3f2}
\caption{Confronto di $k/W_b^2$ a $z=7.5$ m.}
\label{fig:r3f2}
\end{figure}%
\FloatBarrier%

Per lo stesso motivo di prima, abbiamo una discrepanza nei grafici dell'energia cinetica: avendo $C_\mu$ variabile, il modello Realisable limita l'eccessiva produzione di energia turbolenta al centro.

\begin{figure}[H]
\centering
\includegraphics[width = 0.6\textwidth]{r3f3}
\caption{Confronto di $\varepsilon R/W_b^3$ a $z=7.5$ m.}
\label{fig:r3f3}
\end{figure}%
\FloatBarrier%

\newpage

\item Low-Reynolds approach con grid indipendence study

Using the \emph{Null objects} feature, we created three consecutive meshes as follow:
\begin{enumerate}
    \item Mesh 1: $40 \times 200$, $y^+=0.4254$

    \begin{figure}[H]
    \centering
    \includegraphics[width = 0.7\textwidth]{rl3}\\
    \includegraphics[width = 0.3\textwidth]{rlr3}
    \label{fig:rl3}
    \end{figure}

    \item Mesh 2: $50 \times 400$, $y^+=0.271$

    \begin{figure}[H]
    \centering
    \includegraphics[width = 0.7\textwidth]{rl32}
    \\
    \includegraphics[width = 0.3\textwidth]{rlr32}
    \label{fig:rl32}
    \end{figure}

    \item Mesh 3: $60 \times 600$, $y^+=0.231$

    \begin{figure}[H]
    \centering
    \includegraphics[width = 0.7\textwidth]{rl33} \\
    \includegraphics[width = 0.4\textwidth]{rlr33}
    \label{fig:rl33}
    \end{figure}
\end{enumerate}

Per quanto riguarda la convergenza nel caso del modello Low-Reynolds, è importante interpretare correttamente il comportamento dei residui. In particolare, il valore della dissipazione turbolenta \( \varepsilon \) lungo la condotta presenta un \textbf{andamento eterogeneo}. Nella regione completamente sviluppata, \( \varepsilon \) tende a zero all'asse, mostra un picco vicino alle pareti (a causa del viscous sublayer) e risulta molto basso al centro del tubo. Questo comportamento è tipico per un flusso in condotta.

Quando \( \varepsilon \) viene normalizzato rispetto al valore all'ingresso, il trend apparente può risultare amplificato, poiché il valore all'ingresso è molto piccolo, essendo determinato da condizioni iniziali imposte automaticamente dal software o definite dall'utente. Di conseguenza, il \textbf{rapporto di normalizzazione} potrebbe portare a oscillazioni o picchi nei residui, apparentemente significativi ma non necessariamente problematici.

Quindi, se le soluzioni risultano identiche (entro una tolleranza accettabile), è possibile concludere che la soluzione è convergente e valida, anche se i residui mostrano un comportamento oscillatorio.

Questo comportamento dei residui può essere attribuito alla complessità del modello Low-Re rispetto al Wall Function. Nel caso Low-Re, il sistema numerico è più complesso, con celle più piccole vicino al muro e un modello matematico più dettagliato. Di contro, il modello Wall Function, essendo più semplice e con una griglia più grossolana, tende a mostrare residui più stabili, pur sacrificando parte della descrizione fisica del flusso.

Infine, è importante sottolineare che in geometrie complesse e casi reali è normale osservare residui non perfettamente lisci. Quando si raggiunge un \textbf{plateau nel comportamento dei residui}, è spesso sufficiente accettare la soluzione e considerare che il sistema ha raggiunto la convergenza numerica.

\begin{figure}[H]
\centering
\includegraphics[width = 0.5\textwidth]{rgf3}
\caption{Confronto della pressione $P^*$ al centro della tubazione.}
\label{fig:rgf3}
\end{figure}%
\FloatBarrier%

\begin{figure}[H]
\centering
\includegraphics[width = 0.5\textwidth]{rgf2}
\caption{Confronto di $W/W_b$, ascissa in scala logaritmica.}
\label{fig:rgf2}
\end{figure}%
\FloatBarrier%

\begin{figure}[H]
\centering
\includegraphics[width = 0.5\textwidth]{rgf8}
\caption{Confronto di \texttt{DWDZ}.}
\label{fig:rgf8}
\end{figure}%
\FloatBarrier%

\begin{figure}[H]
\centering
\includegraphics[width = 0.6\textwidth]{rgf5}
\caption{Confronto di $w^+$, ascissa in scala logaritmica.}
\label{fig:rgf5}
\end{figure}%
\FloatBarrier%

\begin{figure}[H]
\centering
\includegraphics[width = 0.6\textwidth]{rgf6}
\caption{Confronto di $w^+$, zoom.}
\label{fig:rgf6}
\end{figure}%
\FloatBarrier%

\begin{figure}[H]
\centering
\includegraphics[width = 0.5\textwidth]{rgf7}
\caption{Confronto di $\tau_{\text{visc}}/\tau_w$, $\tau_{\text{Re}}/\tau_w$ e $\tau/\tau_w$.}
\label{fig:rgf7}
\end{figure}%
\FloatBarrier%

Gli sforzi viscosi rimangono a zero al centro, ma si alzano vicino alla parete catturando il comportamento del viscous sublayer. Gli sforzi di Reynolds, associati alle fluttuazioni turbolente, si azzerano al muro, dove il flusso è laminare. Di conseguenza, lo sforzo totale riflette questa combinazione.

\begin{figure}[H]
\centering
\includegraphics[width = 0.5\textwidth]{rgf4}
\caption{Confronto di $f$.}
\label{fig:rgf4}
\end{figure}%
\FloatBarrier%

Anche qui, $f$ diminuisce all'aumentare del numero di elementi.

\begin{figure}[H]
\centering
\includegraphics[width = 0.5\textwidth]{rgf9}
\caption{Confronto di $k/W_b^2$.}
\label{fig:rgf9}
\end{figure}%
\FloatBarrier%

L'energia cinetica turbolenta cresce fino alla zona fuori dal viscous sublayer(per via degli alti gradienti della Reynolds avaraged velocity responsabili della produzione di essa), dove la turbolenza è maggiore, poi va per azzerarsi entrando nella regione laminare vicino al muro.

\begin{figure}[H]
\centering
\includegraphics[width = 0.5\textwidth]{rgf10}
\caption{Confronto di $\varepsilon R/W_b^3$.}
\label{fig:rgf10}
\end{figure}%
\FloatBarrier%

La dissipazione dell'energia cinetica turbolenta cresce vicino al viscous sublayer, dove l'energia viene distrutta, per poi azzerarsi in prossimità del muro, dove il flusso è laminare.

\newpage

\item Confronto delle Mesh 1 tra Wall functions approach e Low-Reynolds approach

\begin{figure}[H]
\centering
\includegraphics[width = 0.5\textwidth]{rcf4}
\caption{Confronto della pressione $P^*$ al centro della tubazione.}
\label{fig:rcf4}
\end{figure}%
\FloatBarrier%

\begin{figure}[H]
\centering
\includegraphics[width = 0.5\textwidth]{rcf1}
\caption{Confronto di $W$ al centro della tubazione.}
\label{fig:rcf1}
\end{figure}%
\FloatBarrier%

\begin{figure}[H]
\centering
\includegraphics[width = 0.5\textwidth]{rcf2}
\caption{Confronto di \texttt{DWDZ}.}
\label{fig:rcf2}
\end{figure}%
\FloatBarrier%

\begin{figure}[H]
\centering
\includegraphics[width = 0.5\textwidth]{rcf3}
\caption{Confronto di $W/W_b$.}
\label{fig:rcf3}
\end{figure}%
\FloatBarrier%

Nel caso Low-Re, il modello cattura meglio il comportamento al muro. Il profilo di velocità risulta più a \emph{panettone} rispetto al Wall Function, poiché vengono considerati i gradienti più ripidi vicino alla parete.

\begin{figure}[H]
\centering
\includegraphics[width = 0.5\textwidth]{rcf6}
\caption{Confronto di $w^+$, ascissa in scala logaritmica.}
\label{fig:rcf6}
\end{figure}%
\FloatBarrier%

\begin{figure}[H]
\centering
\includegraphics[width = 0.5\textwidth]{rcf7}
\caption{Confronto di $\tau_{\text{visc}}/\tau_w$, $\tau_{\text{Re}}/\tau_w$ e $\tau/\tau_w$.}
\label{fig:rcf7}
\end{figure}%
\FloatBarrier%

\begin{figure}[H]
\centering
\includegraphics[width = 0.5\textwidth]{rcf5}
\caption{Confronto di $f$.}
\label{fig:rcf5}
\end{figure}%
\FloatBarrier%

\begin{figure}[H]
\centering
\includegraphics[width = 0.5\textwidth]{rcf8}
\caption{Confronto di $k/W_b^2$.}
\label{fig:rcf8}
\end{figure}%
\FloatBarrier%

\begin{figure}[H]
\centering
\includegraphics[width = 0.5\textwidth]{rcf9}
\caption{Confronto di $\varepsilon R/W_b^3$.}
\label{fig:rcf9}
\end{figure}%
\FloatBarrier%

\end{enumerate}

Possiamo concludere che il Low-Reynolds approach è più accurato, tuttavia bisogna anche considerare i costi computazioni. Se l'obiettivo dell'analisi è valutare gli sforzi ha senso investire, se invece è solo la velocità allora (forse) ha più senso il Wall function approach.



 
%!TEX root = ../main.tex

\section{Quesito 7}

\subsection*{Una domanda con importanti implicazioni pratiche (opzionale)}

\vspace{0.5em}

L'approccio \emph{wall function} impiegato in questo caso di studio richiede che le celle vicine alla parete siano caratterizzate da un $y^+$ superiore a 30 e inferiore a 130. Cosa succede se questa condizione non viene soddisfatta? Si noti che, in pratica, in una geometria complessa questa condizione non sarà soddisfatta per ogni cella.

Eseguire le simulazioni numeriche con il modello di turbolenza $k$-$\varepsilon$ standard ignorando intenzionalmente questo vincolo, ovvero calcolare la soluzione CFD su mesh fini vicino alla parete nel range $y^+ = 1 \div 130$. Analizzare cosa accade quando $y^+$ è al di fuori del range ammissibile (cioè per $y^+$ inferiore a 30), facendo riferimento a tutte le variabili di output.

\NotionDivider

Applicare la legge logaritmica dove vige quella lineare genera turbolenza dove il flusso dovrebbe essere quasi laminare (sottostrato viscoso), con importanti conseguenze come:
\begin{itemize}
    \item senza smorzamento (dumping functions), il codice calcolerà una produzione di turbolenza enorme e irreale attaccata al muro;
    \item friction factor sbagliato;
    \item profili di velocità distorti vicino alla parete;
    \item Convergenza fifficile perché le equazioni sono instabili.
\end{itemize} \cleardoublepage
%!TEX root = ../main.tex

\chapter{Calibration of a load cell}
\thispagestyle{empty}

\section*{Descrizione del caso e obiettivi}

La cella di carico è uno strumento utilizzato per misurare una forza, che produce un'uscita elettrica (voltaggio) approssimativamente proporzionale alla forza applicata (Figura~\ref{fig:c3f1}). Per stimare il valore di una forza tramite una cella di carico, è necessario determinare la \textbf{funzione di trasferimento} che converte il voltaggio in Newton \( V \, [\text{volt}] \rightarrow F \, [\text{Newton}] \). La funzione di trasferimento si ottiene utilizzando campioni di massa il cui peso \( F^* \) è noto con precisione (vedi Tabella~\ref{tab:tab-ref}). Questa operazione viene talvolta definita \textbf{calibrazione} della cella di carico. In effetti, i valori di \( V \) presentano una variabilità casuale, che rende la procedura per determinare la funzione di trasferimento analoga a quella per determinare una \textbf{curva di calibrazione}; ciò spiega il motivo per cui il termine calibrazione viene utilizzato anche in questo contesto.

\begin{figure}[H]
\centering
\includegraphics[width = 0.3\textwidth]{c3f1}
\caption{Rappresentazione schematica.}
\label{fig:c3f1}
\end{figure}

Ogni campione di massa viene applicato alla cella di carico per alcuni secondi e l'uscita elettrica corrispondente viene registrata (la frequenza di campionamento è \( f_s = 100 \, \text{Hz} \)). Questo produce un segnale rumoroso come quello mostrato in Figura~\ref{fig:c3f2} a titolo di esempio.

\begin{figure}[htb]
\centering
\includegraphics[width = 0.5\textwidth]{c3f2}
\caption{Segnale d'esempio.}
\label{fig:c3f2}
\end{figure}

Le storie temporali dell'uscita in tensione e i corrispondenti valori della forza di riferimento \( F^* \) sono elencati nella Tabella~\ref{tab:tab-ref} e forniti sotto forma di \emph{workspace} MATLAB. Ogni workspace contiene il vettore delle letture in Volt (ricordare che la frequenza di campionamento è \( f_s = 100 \, \text{Hz} \)). Nella tabella, il segno di \( F^* \) è rappresentato come mostrato nello schema~\ref{fig:c3f1}: \( F^* > 0 \) quando la forza è diretta verso l'alto, e \( F^* < 0 \) quando la forza è diretta verso il basso.

\begin{table}[H]
    \centering
    \begin{tabular}{|l|c|}
        \toprule
        TestID & $F^*$ [N]  \\
        \midrule
        calibr01 & 0.000    \\
        calibr02 & -0.314   \\
        calibr03 & -0.411   \\
        calibr04 & -0.695   \\
        calibr05 & -0.891   \\
        calibr06 & -1.186   \\
        calibr07 & -1.676  \\
        calibr08 & -2.656   \\
        calibr09 & -5.598   \\
        calibr10 & -10.501  \\
        calibr11 & 0.126    \\
        calibr12 & 1.097    \\
        calibr13 & 2.078    \\
        calibr14 & 4.039   \\
        calibr15 & 8.942   \\
        calibr16 & 13.845  \\
        \bottomrule
    \end{tabular}
    \caption{Dati per la calibrazione.}
    \label{tab:tab-ref}
\end{table}

\section*{Quesito 1}

Determinare la funzione di trasferimento per la cella di carico e stimare l'incertezza dello strumento.

\begin{itemize}
\item Il comportamento rumoroso del segnale di tensione è dovuto alla variabilità dell'uscita elettrica. Il rumore elettrico deve essere filtrato prima di trovare i valori di \( V_0 \) da correlare con \( F^* \) attraverso la funzione di trasferimento. Questo viene ottenuto calcolando i grafici della \textbf{media cumulativa} delle uscite istantanee della tensione.

Tali grafici mostrano, per ogni \( t_i \), la media delle uscite in tensione da 0 a \( t_i \), con il tempo sull'asse orizzontale in scala logaritmica.

Attraverso un'ispezione visiva delle curve, è possibile valutare se il tempo di acquisizione del segnale è sufficientemente lungo da fornire una media cumulativa della tensione \emph{ragionevolmente stabile}. In tal caso, il valore alla fine della curva, \( V \), sarà la tensione \( V_0 \) associata a \( F^* \) nelle successive elaborazioni (in alternativa, si potrebbe considerare una media dei valori nella regione \emph{stabile}). Questa operazione deve essere ripetuta per tutti i test. Una stima approssimativa del tempo di acquisizione richiesto per filtrare il rumore elettrico potrebbe essere ottenuta.

\NotionDivider

Si osservono i seguenti grafici. La media cumulativa può essere considerata stabile dopo un certo intervallo. Tale intervallo, detto $\widehat{t}$, è stato stimato qualitativamente per ogni figura, e il valore finale $V_0$ per ogni test è mostrato nella Tabella~\ref{tab:calibration_results}. In media, abbiamo che:
\begin{equation*}
\widehat{t}_{\text{medio}}=5.25~\text{s}\,.
\end{equation*}

Sempre dal punto di vista qualitativo, si potrebe decidere di non considerare alcuni test perché affetti da \emph{troppo} rumore, per esempio i test 03, 04 e 09 che nella parte \emph{stabile} presentano ancora oscillazioni. Tuttavia, in presenza di \emph{pochi} esperimenti, abbiamo deciso di considerarli tutti.

\begin{figure}[H]
\centering
\includegraphics[width = 0.8\textwidth]{calib-31/test1}
\caption{Media cumulativa Test 01.}
\label{fig:calib31-1}
\end{figure}

\begin{figure}[H]
\centering
\includegraphics[width = 0.8\textwidth]{calib-31/test2}
\caption{Media cumulativa Test 02.}
\label{fig:calib31-2}
\end{figure}

\begin{figure}[H]
\centering
\includegraphics[width = 0.8\textwidth]{calib-31/test3}
\caption{Media cumulativa Test 03.}
\label{fig:calib31-3}
\end{figure}

\begin{figure}[H]
\centering
\includegraphics[width = 0.8\textwidth]{calib-31/test4}
\caption{Media cumulativa Test 03.}
\label{fig:calib31-4}
\end{figure}

\begin{figure}[H]
\centering
\includegraphics[width = 0.8\textwidth]{calib-31/test5}
\caption{Media cumulativa Test 05.}
\label{fig:calib31-5}
\end{figure}

\begin{figure}[H]
\centering
\includegraphics[width = 0.8\textwidth]{calib-31/test6}
\caption{Media cumulativa Test 06.}
\label{fig:calib31-6}
\end{figure}

\begin{figure}[H]
\centering
\includegraphics[width = 0.8\textwidth]{calib-31/test6}
\caption{Media cumulativa Test 06.}
\label{fig:calib31-6}
\end{figure}

\begin{figure}[H]
\centering
\includegraphics[width = 0.8\textwidth]{calib-31/test8}
\caption{Media cumulativa Test 08.}
\label{fig:calib31-8}
\end{figure}

\begin{figure}[H]
\centering
\includegraphics[width = 0.8\textwidth]{calib-31/test9}
\caption{Media cumulativa Test 09.}
\label{fig:calib31-9}
\end{figure}

\begin{figure}[H]
\centering
\includegraphics[width = 0.8\textwidth]{calib-31/test10}
\caption{Media cumulativa Test 10.}
\label{fig:calib31-10}
\end{figure}

\begin{figure}[H]
\centering
\includegraphics[width = 0.8\textwidth]{calib-31/test11}
\caption{Media cumulativa Test 11.}
\label{fig:calib31-11}
\end{figure}

\begin{figure}[H]
\centering
\includegraphics[width = 0.8\textwidth]{calib-31/test12}
\caption{Media cumulativa Test 12.}
\label{fig:calib31-12}
\end{figure}

\begin{figure}[H]
\centering
\includegraphics[width = 0.8\textwidth]{calib-31/test13}
\caption{Media cumulativa Test 13.}
\label{fig:calib31-13}
\end{figure}

\begin{figure}[H]
\centering
\includegraphics[width = 0.8\textwidth]{calib-31/test14}
\caption{Media cumulativa Test 14.}
\label{fig:calib31-14}
\end{figure}

\begin{figure}[H]
\centering
\includegraphics[width = 0.8\textwidth]{calib-31/test15}
\caption{Media cumulativa Test 15.}
\label{fig:calib31-15}
\end{figure}

\begin{figure}[H]
\centering
\includegraphics[width = 0.8\textwidth]{calib-31/test16}
\caption{Media cumulativa Test 16.}
\label{fig:calib31-16}
\end{figure}

\begin{table}[H]
    \centering
    \begin{tabular}{|l|c|c|c|}
        \toprule
        TestID & $F^*$ [N] & $\widehat{t}$ [s] & $V_0$ [V]  \\
        \midrule
        calibr01 & 0.000   & 2  & -0.000027223139966 \\
        calibr02 & -0.314  & 1  & 0.000100843879466 \\
        calibr03 & -0.411  & 1  & 0.000136267348651 \\
        calibr04 & -0.695  & 1  & 0.000248955271896 \\
        calibr05 & -0.891  & 8  & 0.000334117541684 \\
        calibr06 & -1.186  & 6  & 0.000455141343295 \\
        calibr07 & -1.676  & 8  & 0.000656249447185 \\
        calibr08 & -2.656  & 1  & 0.001053949935798 \\
        calibr09 & -5.598  & 1  & 0.002252104292335 \\
        calibr10 & -10.501 & 2  & 0.004253872650178 \\
        calibr11 & 0.126   & 8  & -0.000078179149378 \\
        calibr12 & 1.097   & 2  & -0.000475867746645 \\
        calibr13 & 2.078   & 1  & -0.000874618849803 \\
        calibr14 & 4.039   & 1  & -0.001677096219087 \\
        calibr15 & 8.942   & 35 & -0.003677691934402 \\
        calibr16 & 13.845  & 6  & -0.005671186143606 \\
        \bottomrule
    \end{tabular}
    \caption{Riepilogo dei dati di calibrazione con stima del tempo di acquisizione $\widehat{t}$ e tensione limite $V_0$.}
    \label{tab:calibration_results}
\end{table}

\item Si noti che i valori di \( V_0 \) sono, di per sé, casuali. Infatti, la ripetizione dello stesso test produce valori differenti di \( V_0 \). È possibile derivare una formula che stimi il valore medio limite dei valori di \( V_0 \), indicato come \( \mu_{V_0} \), come funzione del valore di forza di riferimento, \( F^* \). Questa relazione può essere ottenuta tramite regressione lineare delle coppie \( (F^*, V_0) \). La formula avrà la seguente forma:
\begin{equation}
\label{eq:regression}
\mu_{V_0} = m F^* + b
\end{equation}

Determinare i valori di \( m \) e \( b \).

\NotionDivider

Usando il comando \texttt{polyfit} di MATLAB con opzione \emph{degree}$=1$, abbiamo trovato che
\[
m = -0.000407713~\text{V/N} \,, \qquad b = -2.8934 \cdot 10^{-5} ~\text{V}\,.
\]

La retta di regressione così definita è mostrata in Figura~\ref{fig:calib31-rego}.

\begin{figure}[H]
\centering
\includegraphics[width = 0.5\textwidth]{calib-31/rego}
\caption{Retta di regressione lineare.}
\label{fig:calib31-rego}
\end{figure}

\item Tuttavia, l'interesse reale è una formula in cui, dato un valore di tensione \( V_0 \), viene fornita una stima della \emph{vera} forza corrispondente \( F \). Questa formula è chiamata funzione di trasferimento lineare e può essere ottenuta risolvendo la seguente equazione:
\begin{equation}
\label{eq:tildeF1}
\tilde{F} = \frac{V_0 - b}{m}   
\end{equation}

Il simbolo sopra \( F \) indica che, per un dato \( V_0 \), la formula sopra riportata non produce il valore reale di \( F \), ma solo una stima. Questo è dovuto alla natura casuale del processo di misura (si noti che, in linea di principio, lo stesso \( V_0 \) potrebbe corrispondere a diversi possibili valori di \( F \)). 

Ora si desidera calcolare una stima dell'incertezza associata alla stima \( \tilde{F} \). Vale a dire, si desidera trovare \( U \) in modo che, per ogni \( V_0 \), la forza corrispondente rientri nell'intervallo \( \tilde{F} \pm U \).

\begin{enumerate}
\item Un primo approccio consiste nel definire \( U \) come la massima differenza assoluta tra \( F \) e \( F^* \) per tutti i 16 casi di calibrazione. Quale sarebbe il valore di \( U \) in questo caso?

\item Un secondo approccio consiste nel riferirsi agli intervalli di confidenza. Per una distribuzione Gaussiana perfetta, il 95\% delle letture ricade nell'intervallo \( \mu \pm 2\sigma \). Secondo la teoria della regressione lineare, possiamo stimare \( U \) al 95\% di confidenza come:
\begin{equation}
\label{eq:gasu}
2s_f \quad \text{dove} \quad s_f^2 = \frac{1}{N-2} \sum_N \left( \frac{V_0 - b}{m} - F^* \right)^2\,,
\end{equation}

dove \( N \) è il numero totale di campioni.

Quale sarebbe il valore di \( U \) in questo caso? Cosa si potrebbe dedurre confrontando i valori di \( U \) ottenuti con i due metodi?

\emph{\textbf{Nota:} in linea di principio, la formula sopra riportata è ragionevole solo per popolazioni molto ampie, il che non è il caso qui. Per popolazioni moderate, un metodo migliore è disponibile, che si basa sulla distribuzione $t$ di Student. Questo metodo non verrà trattato nel corso base di Fluid Labs.}

\end{enumerate}

\NotionDivider

Nel primo caso calcoliamo $U_1$ come
\begin{equation}
\label{eq:UUU}
U_1=\max \left| F^*- \frac{V_0-b}{m}  \right|=0.01342~\text{N}\,.
\end{equation}

Nel secondo caso, calcoliamo $U_2$ come
\begin{equation}
\label{eq:UUUUU}
U_2=2\sqrt{\frac{1}{N-2} \sum_N \left( \frac{V_0 - b}{m} - F^* \right)^2}=0.01124~\text{N}\,.
\end{equation}

Il primo approccio rappresenta una stima conservativa dell'incertezza, poiché considera la massima deviazione possibile tra \( \tilde{F} \) e \( F^* \). Tuttavia, potrebbe sovrastimare \( U \) in presenza di outlier nei dati.

Invece il secondo approccio, basato sugli intervalli di confidenza, è più robusto poiché tiene conto della distribuzione complessiva dei residui. Tuttavia, non avendo usato un campione ampio non possiamo scartare la prima stima e favorire solamente la seconda.

In ogni caso, abbiamo differenze nell'incertezza di circa \( 0.002 \, \text{N} \) tra i due metodi, considerabile accettabile nel nostro caso.
\end{itemize}

\newpage

\section*{Quesito 2}

Si può stimare la risoluzione della cella di carico? Si ricorda che la risoluzione di uno strumento è la più piccola variazione che può essere misurata. La risoluzione stabilisce i possibili valori di output della cella di carico, con un andamento \emph{a gradini}. Qual è la risoluzione della cella di carico in termini di tensione? E in termini di forza?

\NotionDivider

La risoluzione in tensione è stata stimata analizzando i dati grezzi di \( V \) (tensione) forniti per ciascun test. L'approccio utilizzato si basa sui seguenti passaggi:
\begin{itemize}
    \item I dati grezzi sono stati divisi in finestre temporali di durata pari a 5 secondi.
    \item Per ciascuna finestra, è stata calcolata la media dei valori di tensione.
    \item Le differenze assolute tra le medie delle finestre consecutive sono state calcolate, e il valore minimo tra queste differenze è stato considerato come la risoluzione per il test.
    \item La risoluzione media in tensione è stata calcolata come la media delle risoluzioni ottenute per i 16 test.
\end{itemize}

Il risultato ottenuto per la risoluzione in tensione è:
\[
\Delta V_0 = 6.5512 \cdot 10^{-7}~ \text{V}\,.
\]
La risoluzione in forza è stata calcolata utilizzando la funzione di trasferimento lineare:
\[
\Delta \tilde{F} = \frac{\Delta V_0}{m}=0.00161~\text{N}\,.
\]

Il termine \(b\) non è considerato in quanto non influisce sulle variazioni, essendo una costante.

I risultati ottenuti indicano che la cella di carico ha una risoluzione molto alta in termini di tensione (\( \Delta V_0 \)) e una risoluzione pratica accettabile in termini di forza (\( \Delta \tilde F \)). Questi valori sono in linea con le aspettative per strumenti di misura di alta precisione.




 \cleardoublepage
%!TEX root = ../main.tex

\chapter{Analysis of the velocity field in the wake of a circular cylinder}
\thispagestyle{empty}

\section*{Descrizione e obiettivi del test case}

L'obiettivo di questo laboratorio è analizzare il campo di velocità attorno a un cilindro in un flusso all'interno di un canale d'acqua, concentrandosi sulla \emph{wake region}.

Si fa riferimento ai dati sperimentali ottenuti attraverso la tecnica della \textbf{Particle Streak Velocimetry (PSV)}, che consiste essenzialmente nell'inseminare il flusso con piccole particelle e riprendere un video del piano illuminato da un fascio di luce. La telecamera ha un tempo di esposizione elevato, inteso come la durata temporale di un singolo fotogramma, in modo che ogni particella lasci una scia (\emph{streak}) sull'immagine, come esemplificato in Figura~\ref{fig:test4esempioPSV}. Le immagini vengono poi elaborate per ottenere il campo di velocità.

\begin{figure}[H]
\centering
\includegraphics[width = 0.8\textwidth]{test4img/img1}
\caption{Immagine di esempio dell'esperimento di PSV.}
\label{fig:test4esempioPSV}
\end{figure}

Nel caso di studio proposto, il numero di Reynolds del cilindro $\text{Re}_D=DU_{\infty}/U$ rientra nell'intervallo del \textbf{regime sub-critico}; pertanto, non sarà sorprendente notare che il flusso si separa a una certa distanza dal punto di ristagno anteriore, causando una zona di ricircolo dietro di esso, e che si crea una scia turbolenta oscillante dovuta al distacco di due vortici controrotanti.

Lo schema dell'esperimento è riportato in Figura~\ref{fig:test4sketchPSV}. Il flusso va da destra a sinistra per essere coerente con la posizione del sistema di acquisizione PSV nel setup del canale d'acqua. L'inquadratura della telecamera, che costituisce il campo di misura, è un piano verticale nella sezione mediana del canale. Il sistema di coordinate dell'acquisizione PSV è $(X,Y)$, come indicato nella figura.

\begin{figure}[H]
\centering
\includegraphics[width = \textwidth]{test4img/img2}
\caption{Rappresentazione schematica dell'esperimento.}
\label{fig:test4sketchPSV}
\end{figure}

\section*{Input data}

I dati dell'esperimento sono riportati in Tabella~\ref{tab:tab4-dati}.

\begin{table}[H]
    \centering
    \begin{tabular}{clcc}
        \toprule
        \textbf{Symbol} & \textbf{Parameter} & \textbf{Value} & \textbf{Units}  \\
        \midrule
        $B$ & Width of the chanell & 0.5 & m    \\
        $D$ & Diameter of the cylinder & 0.06 & m \\
        $f_s$ & Sampling frequency & 50 & Hz \\
        $h$ & Water level in the channel upstream of the cylinder & 0.42 & m \\
        $h_b$ & Distance of the cylinder wall from the channel bottom & 0.18 & m \\
        $Q$ & Volumetric flow rate of water & 35 & l/s \\
        res & Resolution of the images & 3040 & px/m \\
        $\rho$ & Density of water & 998 & kg/m$^3$  \\
        $\mu$ & Dynamic viscosity of water & 0.001 & Pa$\cdot$s \\
        \bottomrule
    \end{tabular}
    \caption{Dati dell'esperimento.}
    \label{tab:tab4-dati}
\end{table}

Per la post-elaborazione dei dati PSV, l'inquadratura è stata suddivisa in una griglia regolare di $30 \times 23$ celle uniformi lungo le direzioni $X$ e $Y$. Il numero di passi temporali è 1499, e la durata di ciascuno di essi è pari a $1/f_s$. Il set completo dei dati di acquisizione è fornito nel workspace MATLAB \texttt{PSVdata.mat}, che include le variabili mostrate in Tabella~\ref{tab:tab4-PSV}.

\begin{table}[h!]
\centering
\begin{tabular}{lll}
\toprule
\textbf{Variabile} & \textbf{Tipo} & \textbf{Valore} \\
\midrule
\texttt{Grid\_Xpx} & matrice $23 \times 30$ & Coordinate $X$ dei centri della griglia [in px] \\
\texttt{Grid\_Ypx} & matrice $23 \times 30$ & Coordinate $Y$ dei centri della griglia [in px] \\
\texttt{SXpx} & matrice $23 \times 30 \times 1499$ & Spostamento $X$ [in px] per ogni passo temporale \\
\texttt{SYpx} & matrice $23 \times 30 \times 1499$ & Spostamento $Y$ [in px] per ogni passo temporale \\
\bottomrule
\end{tabular}
\caption{Dati PSV.}
\label{tab:tab4-PSV}
\end{table}

Un'istantanea statica dell'inquadratura PSV con un righello vicino al cilindro è stata scattata prima dell'esecuzione dell'esperimento, ed è fornita come file immagine \texttt{ref\_image.png}. Una piccola anteprima del file immagine è riportata in Figura~\ref{fig:test4refimg}. Lo scopo di questo passaggio preliminare è duplice. Da un lato, è stato utilizzato per stimare la risoluzione dell'immagine come $\text{res} = 3040~\text{px/m}$, che è qui già fornita come input per semplicità. Dall'altro lato, serve per trasformare le matrici \emph{pixel-based} \texttt{Grid\_Xpx}, \texttt{Grid\_Ypx}, \texttt{SXpx} e \texttt{SYpx} nelle corrispondenti matrici dimensionali nel nuovo sistema di coordinate ($x_0, y_0$), centrato nel punto di ristagno posteriore e orientato come mostrato in Figura~\ref{fig:test4sketchPSV}.

\begin{figure}[H]
\centering
\includegraphics[width = 0.5\textwidth]{test4img/ref_image}
\caption{Anteprima dell'immagine di riferimento.}
\label{fig:test4refimg}
\end{figure}

\section*{Analisi preliminare}

\subsection*{Quesito 1}

Calcolare la velocità media del canale (bulk velocity) $U_b$ basandosi sulle quantità misurate a monte del cilindro $Q$, $B$, $h$. Assumere che questa velocità sia quella della corrente libera che approccia il cilindro, $U_{\infty}$, e poi calcolare il numero di Reynolds $Re_D = D \cdot U_{\infty} / \nu$. Notare che $U_b \neq U_{\infty}$ poiché gli esperimenti coinvolgono un flusso d'acqua in un canale di dimensioni finite. Ciononostante, questa potrebbe essere una ragionevole prima approssimazione.

\NotionDivider

La velocità media del flusso nel canale è calcolata come:
\[
U_b = \frac{Q}{B \cdot h} = \frac{35\cdot 10^{-3}}{0.5\cdot 0.42}  = 0.1667 ~ \text{m/s}\,.
\]

Il numero di Reynolds è:
\[
\text{Re}_D = \frac{D \cdot U_b \cdot \rho}{\mu} = \frac{0.06 \cdot 0.167 \cdot 998}{0.001} = 9980\,.
\]

Poiché il numero di Reynolds  rientra nell'intervallo $5 \cdot 10^3 < \text{Re}_D < 3 \cdot 10^5$, il regime del flusso può essere classificato come subcritico.

\subsection*{Quesito 2}

Trasformare le matrici \emph{pixel-based} \texttt{Grid\_Xpx}, \texttt{Grid\_Ypx}, \texttt{SXpx} e \texttt{SYpx} in appropriate matrici adimensionali. Si fa ora riferimento al sistema di coordinate ($x_0, y_0$), che è più fisico rispetto a quello del frame PSV ($X, Y$). A tal fine, seguire i seguenti passaggi:

\begin{enumerate}
\item Trovare le coordinate dell'origine del nuovo sistema di coordinate ($x_0, y_0$), sia in pixel che in metri applicando la risoluzione suggerita $\text{res} = 3040 \text{ px/m}$.

\item Partire da \texttt{Grid\_Xpx} e \texttt{Grid\_Ypx} e costruire due nuove matrici \texttt{Grid\_x0m}, \texttt{Grid\_y0m}, che contengono il centro delle celle espresso in metri secondo il nuovo sistema di coordinate ($x_0, y_0$). Questo richiede la conversione delle vecchie matrici \emph{pixel-based} in matrici dimensionali (in metri), e poi l'applicazione della trasformazione del sistema di coordinate per passare da ($X, Y$) a ($x_0, y_0$) --- non necessariamente in questo ordine. Infine, costruire le matrici adimensionali \texttt{Grid\_x0\_D} e \texttt{Grid\_y0\_D} ottenute normalizzando le coordinate dei punti della griglia con il diametro del cilindro, $D$.

\item Partire da \texttt{SXpx} e \texttt{SYpx} e costruire nuove matrici \texttt{Sx0m} e \texttt{Sy0m}, contenenti gli spostamenti in metri secondo il nuovo sistema di coordinate ($x_0, y_0$). Ottenere le matrici del campo di velocità \emph{instantaneo} $u_0$, $v_0$ considerando che gli spostamenti sono stati misurati durante il tempo di esposizione $1/f_s$. Infine, rendere le due matrici adimensionali dividendole per la velocità media $U_b$; chiamiamo le matrici finali \texttt{u0\_Ub}, \texttt{v0\_Ub}.
\end{enumerate}

\NotionDivider

Il punto \((x_0, y_0)\) nella griglia corrisponde al pixel (465, 285), come è evidenziato in Figura~\ref{fig:test4refpxl}. Convertendo queste coordinate in metri, si ottiene:
\[
x_0 = \frac{465}{3040} = 0.1530 ~ \text{m}\,, \quad y_0 = \frac{285}{3040} = 0.0938 ~ \text{m}\,.
\]

\begin{figure}[H]
\centering
\includegraphics[width = 0.5\textwidth]{test4img/ref_pxl}
\caption{Immagine di riferimento con centro $(x_0, y_0)$ evidenziato.}
\label{fig:test4refpxl}
\end{figure}

Per centrare il sistema di coordinate nel punto \((x_0, y_0)\), che rappresenta il punto di stagnazione posteriore del cilindro, le coordinate della griglia vengono trasformate come segue:
\[
\texttt{Grid\_x0m} = -\left( \frac{\texttt{Grid\_Xpx}}{3040} - 0.1530 \right), \quad
\texttt{Grid\_y0m} = -\left( \frac{\texttt{Grid\_Ypx}}{3040} - 0.0938 \right).
\]

Le coordinate sono successivamente normalizzate rispetto al diametro del cilindro \(D\):
\[
\texttt{Grid\_x0\_D} = \frac{\texttt{Grid\_x0m}}{D}\,, \quad \texttt{Grid\_y0\_D} = \frac{\texttt{Grid\_y0m}}{D}\,.
\]

I vettori di spostamento inizialmente in pixel, vengono convertiti in metri:
\[
\texttt{Sx0m} = -\frac{\texttt{SXpx}}{3040}\,, \quad \texttt{Sy0m} = -\frac{\texttt{SYpx}}{3040}\,.
\]

\emph{Nota:} lo shift non va applicato anche qui, essendo che i vettori spostamento esprimono una differenza tra posizioni.

Le velocità istantanee (\(u_0\) e \(v_0\)) vengono calcolate considerando il tempo di esposizione \(1/f_s\):
\[
u_0 = \texttt{Sx0m} \cdot f_s\,, \quad v_0 = \texttt{Sy0m} \cdot f_s\,.
\]

Infine, vengono normalizzate come segue:
\begin{equation*}
\texttt{u0\_Ub} = \frac{u_0}{U_b}\,,\quad  \texttt{v0\_Ub} = \frac{v_0}{U_b}\,.
\end{equation*}




\section*{Analisi del campo velocità media di Reynolds}

\subsection*{Quesito 3}

Calcolare il campo di velocità medio di Reynolds tramite le matrici \texttt{u0\_Ub} e \texttt{v0\_Ub}. Questo richiede di eseguire la media temporale delle matrici di velocità \emph{istantanee} precedentemente calcolate. Si raccomanda di utilizzare il comando MATLAB \texttt{nanmean}, in modo che le numerose celle \texttt{Nan} vengano ignorate. Infine, visualizzare il campo di velocità medio di Reynolds mostrando le seguenti quantità:

\begin{enumerate}
\item Il color plot del modulo della velocità media di Reynolds. A tal fine, si può utilizzare, ad esempio, il comando \texttt{pcolor} con \texttt{shading interp}.

\item La mappa vettoriale tramite il comando MATLAB \texttt{quiver}. Si raccomanda di disegnare anche il cilindro tramite il comando MATLAB \texttt{rectangle}. Notare che si troveranno vettori spuri anche all'interno dello spazio occupato dal cilindro. Chiaramente, questi vettori sono non fisici e sono la conseguenza di errori di prospettiva o problemi nell'algoritmo PSV. Devono essere ignorati.

\item Le streamlines con frecce tramite il comando MATLAB \texttt{streamslice}. Ancora una volta, si raccomanda di disegnare anche il cilindro tramite il comando MATLAB \texttt{rectangle}.

\item I profili di $U_0/U_b$ lungo la linea orizzontale $Y_0/D=0$ e lungo le due linee verticali a $X_0/D=0.5$ e $X_0/D=1.5$. Questo potrebbe richiedere l'interpolazione dei valori della griglia in posizioni spaziali selezionate. Potreste usare il comando MATLAB \texttt{interp2}.
\end{enumerate}

\NotionDivider

Le velocità medie di Reynolds adimensionali sono calcolate come:
\[
\frac{U_0}{U_b} = \texttt{nanmean}(\texttt{u0\_Ub})\,, \quad \frac{V_0}{U_b} = \texttt{nanmean}(\texttt{v0\_Ub})\,.
\]

\emph{Nota:} in realtà usiamo il comando \texttt{U0\_Ub=mean(u0\_Ub,3,'omitnan')}.

A questo punto calcoliamo
\begin{equation*}
\texttt{Velocity\_Magntidue} = \sqrt{\left( \frac{U_0}{U_b} \right)^2 + \left( \frac{V_0}{U_b} \right)^2}
\end{equation*}

e osserviamo la Figura~\ref{fig:test431}.

\begin{figure}[H]
\centering
\includegraphics[width = 0.7\textwidth]{test4img/ravg1}
\caption{Colorplot tramite \texttt{pcolor}.}
\label{fig:test431}
\end{figure}

Nella scia del cilindro osserviamo una regione turbolenta. La distribuzione della magnitudine della velocità evidenzia una diminuzione significativa della velocità nella regione centrale della wake, rappresentata dalle tonalità di colore blu, che corrispondono a velocità più basse. Questo comportamento è tipico della zona di ricircolo dietro il cilindro, dove il flusso inverte localmente la sua direzione. Nelle aree più lontane dal cilindro, le tonalità di colore giallo e verde indicano un graduale recupero della velocità del flusso man mano che la turbolenza si dissipa. Il boundary layer è laminare nella regione iniziale, ma la wake diventa completamente turbolenta, coerente con il numero di Reynolds elevato.

Ora si osservi la Figura~\ref{fig:test432}. La mappa vettoriale e le streamline evidenziano una simmetria quasi perfetta rispetto alla linea \(y_0/D = 0\), come atteso nel regime subcritico. 
Il punto di separazione del boundary layer si verifica teoricamente in prossimità di un angolo di \textbf{82° nella parte frontale del cilindro}. Tuttavia, nella realtà, il punto di separazione può oscillare a causa della turbolenza nella wake. 
La separazione del flusso avviene rapidamente in presenza di un \textbf{gradiente di pressione avversa}, infatti il minimo per la pressione si raggiunge ad 80°, ma il boundary layer laminare è talmente debole da far si che la separazione avvenga subito. Le streamline mostrano chiaramente la formazione di \textbf{due zone di ricircolo} opposte nella scia dietro il cilindro, con la wake caratterizzata da turbolenza altamente irregolare e oscillante.

\begin{figure}[H]
\centering
\includegraphics[width = 0.45\textwidth]{test4img/ravg2} \quad
\includegraphics[width = 0.45\textwidth]{test4img/ravg3}
\caption{Vector map tramite \texttt{quiver} (sinistra). Streamlines tramite \texttt{streamslice} (destra).}
\label{fig:test432}
\end{figure}

Infine, si osservino i grafici di Figura~\ref{fig:test433}. Per quanto riguarda il profilo lungo l'asse orizzontale, si osserva un \textbf{cambiamento di direzione della velocità}, indicativo della transizione verso una wake completamente turbolenta. Le due zone di ricircolo influenzano significativamente il profilo della velocità. Invece, i profili verticali mostrano che, dopo il cilindro, la velocità cambia segno in corrispondenza di \(y_0/D = 0\), segnalando l'ingresso in una zona di turbolenza e ricircolazione. Man mano che ci si allontana dal cilindro (\(x_0/D = 1.5\)), la variazione della velocità è meno drastica, un segno che l'intensità delle turbolenze diminuisce.

\begin{figure}[H]
\centering
\includegraphics[width = 0.45\textwidth]{test4img/ravg4} \quad
\includegraphics[width = 0.45\textwidth]{test4img/ravg5}
\caption{Profilo orizzontale lungo $y_0/D=0$ (sinistra). Profili verticali lungo $x_0/D=0.5$ e $x_0/D=1.5$ (destra).}
\label{fig:test433}
\end{figure}

\section*{Analisi dell'evoluzione dinamica del flusso}

\subsection*{Quesito 4}

Investigare l'evoluzione dinamica del flusso, come segue:

\begin{enumerate}
\item Plottare la storia temporale della velocità verticale normalizzata $v_0/U_b$ in un punto attorno alla parte posteriore del cilindro.

\item Utilizzare la funzione MATLAB \texttt{fft\_of\_v0\_Ub\_velocity} fornita per calcolare la FFT (Fast Fourier Transform) del segnale di velocità $v_0/U_b$ e trovare lo spettro di frequenza. Identificare la frequenza di \emph{vortex shedding}, $f_{\text{shedding}}$, come il picco principale dello spettro di frequenza.

\item Stimare il numero di Strouhal caratteristico della scia oscillante, come $\text{Sr} = f_{\text{shedding}} \cdot D / U_b$, e confrontarlo con il valore di riferimento della letteratura allo stesso numero di Reynolds per il caso non confinato (unbounded).
\end{enumerate}

\NotionDivider

La serie temporale della componente normalizzata $v_0/U_b$ è stata analizzata nel punto $x_0/D \approx 1.5$, $y_0/D = 0$, quindi all'interno della scia, ed è mostrata a sinistra in Figura~\ref{fig:test4fuorier}. L'andamento temporale della velocità istantanea non è stazionario (in senso statistico) quindi si indagano le frequenze del segnale in cerca della frequenza dominante di \textbf{vortex shedding}. Nella figura di destra viene illustrato lo spettro di frequenza calcolato utilizzando la FFT. La frequenza dominante, indicata nel grafico, risultata pari a $f_{\text{shedding}} = 0.57$ Hz.

\begin{figure}[H]
\centering
\includegraphics[width = 0.45\textwidth]{test4img/fuorier1} \quad
\includegraphics[width = 0.45\textwidth]{test4img/fuorier2}
\caption{Storia temporale di $v_0/U_b$, zoom (sinistra). Spettro della frequenza (destra).}
\label{fig:test4fuorier}
\end{figure}

Possiamo infine validare i nostri risultati, confrontandoli con quelli dell'articolo Fey et al. (1998) (\href{https://doi.org/10.1063/1.869675}{doi}).

Utilizzando la relazione empirica per il regime subcritico, il numero di Strouhal teorico per il nostro caso può essere calcolato come:
\[
\text{Sr} = \frac{f_{\text{shedding}} \cdot D}{U_b}
\approx 0.2042\,.
\]

I risultati di Fey et al. mostrano che
\begin{equation*}
\text{Sr}_{\text{Fey}} = \text{Sr}^* + \frac{m}{\sqrt{\text{Re}}}=0.1996\,,\text{ con }\text{Sr}^*=0.1776 \text{ e }m=2.2023\,.
\end{equation*}

\begin{figure}[H]
\centering
\includegraphics[width = 0.7\textwidth]{test4img/fuorier3}
\caption{Paragone tra risultati sperimentali.}
\label{fig:test4fuorier2}
\end{figure}

Osserviamo Figura~\ref{fig:test4fuorier2}. I risultati differiscono del 2.32\%. Tale errore, considerabile comunque accettabile, è dovuto a queste principali differenze:
\begin{itemize}
    \item \textbf{Dominio}: Nel nostro esperimento il flusso avviene in una condotta d'acqua, mentre i dati sperimentali sono stati ottenuti in galleria del vento. La condotta introduce l'effetto dei bordi del dominio, che non è presente in un wind tunnel.
    \item \textbf{Errori Numerici}: I calcoli numerici introducono una discreta incertezza nei risultati.
\end{itemize}

Confermiamo quindi la consistenza fisica dei nostri risultati.







 \cleardoublepage
%!TEX root = ../main.tex

\chapter{Analysis of the hydrodynamic forces acting on a cylinder in steady free surface flow}
\thispagestyle{empty}

\section*{Descrizione e obiettivi del test case}

L'obiettivo di questo caso di studio è investigare le forze di resistenza (\emph{drag}) e portanza (\emph{lift}) agenti su un cilindro sommerso in un flusso d'acqua stazionario a superficie libera. Gli esperimenti sono stati eseguiti nel canale d'acqua del Laboratorio di Idraulica, e le due componenti della forza (orizzontale e verticale) sono state misurate tramite una bilancia. Nel caso di studio proposto, il numero di Reynolds del cilindro $\text{Re}=D U_{\infty}/\nu$ rientra nell'intervallo del \emph{sub-critical regime}; pertanto, non sarà sorprendente notare che il flusso si separa a una certa distanza dal punto di ristagno anteriore, causando una zona di ricircolo dietro di esso, e che una scia oscillante viene creata dallo \emph{shedding} di due vortici controrotanti. Di conseguenza, anche la resistenza e la portanza mostreranno un comportamento oscillatorio.

\begin{figure}[H]
\centering
\includegraphics[width = \textwidth]{test5img/wrek}
\caption{Sketch of the experiment.}
\label{fig:test5sketch}
\end{figure}

Lo schema dell'esperimento è riportato in Figura~\ref{fig:test5sketch}. La bilancia è fissata al suolo e collegata al cilindro tramite le \emph{endplates}. Le \emph{endplates}, sagomate idrodinamicamente, permettono di sostenere il cilindro e assicurare una sorta di \emph{two-dimensional flow} attorno ad esso sopprimendo gli effetti tridimensionali alle estremità del cilindro. Due celle di carico sono installate all'interno della bilancia, e forniscono un voltaggio in uscita linearmente proporzionale alla forza applicata.

La calibrazione della bilancia è stata eseguita dopo l'installazione delle \emph{endplates} e del cilindro sulla bilancia, applicando pesi campione al cilindro senza acqua nel canale come mostrato in Figura~\ref{fig:bello3}. Nello sviluppo della funzione di calibrazione, si è tenuto conto del fatto che le forze reali subite dalle celle non sono solo quelle prodotte dal carico applicato, ma includono anche il peso della struttura e altri piccoli contributi relativi alla deformabilità della struttura stessa. La funzione di calibrazione è stata determinata in modo tale che la condizione $F_x=0$, $F_y=0$ corrisponda all'assenza di forze (esterne) applicate, al netto del peso e dei piccoli contributi legati alla deformabilità.

\begin{figure}[H]
\centering
\includegraphics[width = 0.5 \textwidth]{test5img/bello3}
\caption{Calibration setup.}
\label{fig:bello3}
\end{figure}

Quando l'acqua scorre nel canale, assumendo che le forze dinamiche sulle \emph{endplates} siano trascurabili dato che sono sagomate idrodinamicamente, la forza esterna orizzontale $F_x$ è uguale alla forza di resistenza (\emph{drag}) agente sul cilindro, $F_D$.
Viceversa, la forza esterna verticale $F_y$ è uguale alla somma della forza di portanza (\emph{lift}) agente sul cilindro $F_L$ e la forza di galleggiamento (\emph{buoyancy force}) agente sul cilindro e sulle \emph{endplates}, $F_B$.
Quindi, mentre $F_D$ è semplicemente presa come l'output della calibrazione $F_x$, $F_L$ sarà data da:
\begin{equation*}
F_L = F_y - F_B
\end{equation*}
La forza di galleggiamento, $F_B$, potrebbe essere teoricamente calcolata moltiplicando il volume delle parti immerse (cilindro e parte delle \emph{endplates}) per il peso specifico dell'acqua. Tuttavia, poiché conoscere tutti i dettagli geometrici con alta accuratezza non è banale, misurare direttamente $F_B$ con la bilancia appare un'opzione preferibile. Questo si ottiene effettuando un test in acqua ferma (\emph{still water}) con lo stesso livello del test con acqua in movimento. Poiché nessuna forza di resistenza e portanza gioca un ruolo nel test statico, in questo caso la forza esterna orizzontale $F_x$ sarà zero e la forza esterna verticale $F_y$ sarà uguale a $F_B$.

\newpage

\section*{Dati di Input}

I dati di input del problema sono riassunti nella Tabella~\ref{tab:tab5-dati}.

\begin{table}[H]
    \centering
    \begin{tabular}{clcc}
        \toprule
        \textbf{Symbol} & \textbf{Parameter} & \textbf{Value} & \textbf{Units}  \\
        \midrule
        $B$ & Width of the channel & 0.5 & m    \\
        $D$ & Diameter of the cylinder & 0.06 & m \\
        $L$ & Width of the cylinder & 0.185 & m \\
        $t$ & Thickness of the endplates & 0.002 & m \\
        $b$ & Length of the endplates & 0.180 & m \\
        $h$ & Water level in the channel upstream of the cylinder & 0.45 & m \\
        $h_b$ & Distance of the cylinder wall from the channel bottom & 0.18 & m \\
        $f_s$ & Sampling frequency & 200 & Hz \\
        $Q$ & Volumetric flow rate of water & 75 & l/s \\
        $\rho$ & Density of water & 998 & kg/m$^3$  \\
        $\mu$ & Dynamic viscosity of water & 0.001 & Pa$\cdot$s \\
        \bottomrule
    \end{tabular}
    \caption{Dati dell'esperimento.}
    \label{tab:tab5-dati}
\end{table}

Il set completo dei dati di acquisizione è fornito nei seguenti \emph{workspaces} MATLAB (Tabella~\ref{tab:tab5-work}). In ogni \emph{workspace}, i valori di $F_x$ e $F_y$ sono forniti sotto forma di vettori. Questi valori sono già stati convertiti dal voltaggio in uscita delle celle tramite le funzioni di calibrazione, come spiegato precedentemente.

\begin{table}[H]
    \centering
    \begin{tabular}{ll}
        \toprule
        \textbf{Filename} & \textbf{Condition} \\
        \midrule
        \texttt{FORCEdata\_stillwater.mat} & Cylinder and endplates, still water in the channel \\
        \texttt{FORCEdata\_flow.mat} & Cylinder and endplates, flowing water in the channel \\
        \bottomrule
    \end{tabular}
    \caption{File dei dati sperimentali.}
    \label{tab:tab5-work}
\end{table}

\section*{Quesito 1}

Calcolare la velocità media del canale (\emph{bulk velocity}) $U_b$ basandosi sulle quantità misurate a monte del cilindro $Q$, $B$, $h$. Assumere che questa velocità sia quella della corrente libera che approccia il cilindro, $U_{\infty}$, e poi calcolare il numero di Reynolds $\text{Re}=D U_{\infty}/\nu$.

\NotionDivider

La velocità media del flusso nel canale è calcolata come:
\[
U_b =  \frac{Q}{B \cdot h}  \approx 0.333 ~ \text{m/s}\,.
\]

Il numero di Reynolds risulta quindi:
\[
\text{Re} = \frac{D \cdot U_\infty \cdot \rho}{\mu}  \approx \frac{D \cdot U_b \cdot \rho}{\mu}  \approx 19960\,.
\]
\noindent
Il regime subcritico è confermato poiché \(5 \cdot 10^3 < \text{Re} < 3 \cdot 10^5\).

\newpage

\section*{Quesito 2}

Calcolare la forza di galleggiamento (\emph{buoyancy force}) $F_B$ agente sulle parti immerse (cilindro e parte delle \emph{endplates}).

\NotionDivider

La forza di galleggiamento è determinata utilizzando il test in acqua ferma, denominato \texttt{FORCEdata\_stillwater.mat}. Come mostrato in Figura~\ref{fig:bello2}, abbiamo che:
\[
\boxed{
F_y = F_B \qquad \left( F_x=0 \right)
}
\]

Dai dati del test statico (facciamo la media temporale per filtrare il rumore):
\[
F_B = \frac{1}{12800} \sum_{i=1}^{12800} F_{y,i} = 7.3236 ~ \text{N}\,.
\]

\begin{figure}[H]
\centering
\includegraphics[width = 0.7 \textwidth]{test5img/bello2}
\caption{Still water test.}
\label{fig:bello2}
\end{figure}

\newpage

\section*{Quesito 3}

Calcolare le forze di resistenza (\emph{drag}) e portanza (\emph{lift}) medie di Reynolds agenti sul cilindro e mostrate in Figura~\ref{fig:drl}, $\langle F_D\rangle$ e $\langle F_L\rangle$, e i coefficienti di resistenza e portanza medi di Reynolds, $\langle C_D\rangle$ e $\langle C_L\rangle$. Notare che \textbf{le forze di resistenza e portanza agenti sulle \emph{endplates} sono considerate trascurabili}.

\begin{figure}[H]
\centering
\includegraphics[width = 0.3 \textwidth]{test5img/drl}
\caption{Drag and Lift forces.}
\label{fig:drl}
\end{figure}

\NotionDivider

Ora consideriamo il test in acqua corrente \texttt{FORCEdata\_flow.mat}. Come mostrato in Figura~\ref{fig:bello1}, le forze di resistenza e portanza medie di Reynolds sono
\begin{equation*}
\boxed{
    F_x \approx F_D \qquad F_y \approx F_B+F_L
}  
\end{equation*}

In Matlab, le calcoliamo come:
\begin{equation*}
\begin{aligned}
\langle F_D\rangle &= \frac{1}{62800} \sum_{i=1}^{62800} F_{x,i} = 0.9223 ~ \text{N}\,, \\
\langle F_L\rangle &= \frac{1}{62800} \sum_{i=1}^{62800} F_{y,i}  -F_B=-0.0172 ~ \text{N} \,.
\end{aligned}
\end{equation*}

\begin{figure}[H]
\centering
\includegraphics[width = 0.9 \textwidth]{test5img/bello1}
\caption{Flowing water test.}
\label{fig:bello1}
\end{figure}

I coefficienti di resistenza e portanza medi di Reynolds sono:
\[
\langle C_D \rangle = \frac{\langle F_D \rangle}{\frac{1}{2} \rho U_b^2 A}\,, \qquad \langle C_L \rangle = \frac{\langle F_L \rangle}{\frac{1}{2} \rho U_b^2 A}\,,
\]

dove l'area frontale di riferimento è 
\[
A = D \cdot L = 0.06 \cdot 0.185 = 0.0111 ~ \text{m}^2\,.
\]

I risultati ottenuti sono:
\[
\langle C_D \rangle = 1.4987\,, \qquad \langle C_L \rangle = -0.0280\,.
\]

Coerentemete alle aspettative, poiché il cilindro è un corpo simmetrico e il flusso è mediamente simmetrico, $\langle F_L \rangle$ e $\langle C_L \rangle$ sono molto vicini a 0.

\section*{Quesito 4}

Calcolare le incertezze di $C_D$ e $C_L$ partendo dalla definizione di queste quantità e applicando la legge di propagazione degli errori:
\begin{equation}\label{eq:test5q4}
y=y(x_1,\dots,x_n)\quad\rightarrow\quad
u(y) = \sqrt{\sum_{i=1}^{n} \left( \frac{\partial y}{\partial x_i} u(x_i) \right)^2}
\end{equation}

Fare riferimento ai valori di incertezza nella Tabella~\ref{tab:tab5-un} e provare a stimare i valori mancanti.

\begin{table}[H]
    \centering
    \begin{tabular}{llcc}
        \toprule
        \textbf{Variable} & \textbf{Instrument} & \textbf{Uncertainty} & \textbf{Units} \\
        \midrule
        $B$ & Ruler & ??? & mm \\
        $D$ & Vernier caliber with scale value 0.05 & ??? & mm \\
        $L$ & Ruler & ??? & mm \\
        $h$ & Piezometer & 1 & mm \\
        $Q$ & \parbox[t]{5cm}{Proline Promag 50L installed \\ in a 200 mm pipe (DN200)} & ??? & of reading \\
        $\rho$ & -- & $\approx 0$ &  \\
        $F_x$ & From calibration of frame & 0.045 & N \\
        $F_y$ & From calibration of frame & 0.025 & N \\
        \bottomrule
    \end{tabular}
    \caption{Incertezze strumentali.}
    \label{tab:tab5-un}
\end{table}

\emph{\textbf{Nota:} Si noti che queste formule potrebbero essere applicate a livello istantaneo, ovvero ottenendo l'incertezza per ogni singolo valore $C_D(t)$ e $C_L(t)$; in questo caso, le incertezze dei due coefficienti varieranno nel tempo. Tuttavia, appare più pratico calcolare direttamente le incertezze dei coefficienti mediati alla Reynolds, $\langle C_D\rangle$ e $\langle C_L\rangle$, partendo da $\langle F_D\rangle$ e $\langle F_L\rangle$. Le incertezze per $F_x$ e $F_y$ valgono ragionevolmente anche per i loro valori mediati alla Reynolds, $\langle F_x\rangle$ e $\langle F_y\rangle$.}

\NotionDivider

Per calcolare le incertezze di $\langle C_D\rangle$ e $\langle C_L\rangle$ dobbiamo esprimere tali quantità in funzione di parametri di cui sappiamo l'incertezza. Di conseguenza, abbiamo:
\begin{equation*}
\langle C_D \rangle = \frac{\langle F_D \rangle}{\frac{1}{2} \rho U_b^2 A} = \frac{2 \langle F_D \rangle B^2 h^2}{ \rho Q^2 DL}
\end{equation*}

Le derivate parziali di $\langle C_D \rangle$ rispetto alle variabili sono:
\begin{align*}
\frac{\partial \langle C_D \rangle}{\partial \langle F_D \rangle} &= \frac{2 B^2h^2}{\rho Q^2 A}\,, &\frac{\partial \langle C_D \rangle}{\partial Q} = -\frac{4 \langle F_D \rangle B^2h^2}{\rho Q^3 A}\,, \\
\frac{\partial \langle C_D \rangle}{\partial B} &= \frac{4 \langle F_D \rangle B h^2}{\rho Q^2 A}\,, &\frac{\partial \langle C_D \rangle}{\partial D} = -\frac{2 \langle F_D \rangle B^2h^2}{\rho Q^2 D^2 L}\,, \\
\frac{\partial \langle C_D \rangle}{\partial h} &= \frac{4 \langle F_D \rangle B^2 h}{\rho Q^2 A}\,, &\frac{\partial \langle C_D \rangle}{\partial L} = -\frac{2 \langle F_D \rangle B^2h^2}{\rho Q^2 D L^2}\,.
\end{align*}

\emph{Nota:} non calcoliamo la derivata rispetto $\rho$ perché l'incertezza di $\rho$ è praticamente nulla.

Prima di applicare la formula di propagazione dell'incertezza, dobbiamo trovare le incertezze di tali variabili:
\begin{itemize}
\item $u(\langle F_D \rangle)=u(F_D)=u(F_x)=0.045~$N (Tabella~\ref{tab:tab5-un}).
\item $u(B)=u(L)=\frac{1 ~ \text{mm}}{2} = 0.0005 ~ \text{m}$ (Ikea ruler).
\item $u(D) = 50~\mu\text{m} = 0.00005 ~ \text{m}$ (pagina 1 del file \texttt{caliber-accuracy.pdf})
\item $u(Q)=A_\text{pipe}\cdot u(U_\text{pipe})$, dove $u(U_\text{pipe})$ rappresenta l'incertezza sulla velocità media e, come specificato a pagina 10 del manuale \texttt{proline-promag-50L-datasheet.pdf}, si esprime come
\begin{equation*}
u(U_\text{pipe})= \frac{0.5}{100} \cdot \frac{Q}{A_\text{pipe}} + 0.001 = \frac{0.5}{100} \cdot \frac{Q}{\frac{\pi}{4} D_N^2} + 0.001 \approx 0.0129  ~ \text{m/s}  
\end{equation*}

Dunque, abbiamo che $u(Q)\approx 0.000406~\text{m}^3/\text{s}$.
\end{itemize}

A questo punto, possiamo applicare la formula di propagazione dell'errore~\eqref{eq:test5q4} e otteniamo
\begin{equation*}
u(\langle C_D \rangle)=0.0754 \qquad \Longrightarrow \qquad \langle C_D \rangle = 1.4987 \pm 0.0754\,.
\end{equation*}

Per quanto riguarda $u(\langle C_L \rangle)$ il calcolo è analogo, con l'accortezza che $u(\langle F_L \rangle)=u(F_L)=u(F_B-F_y)=\sqrt{2u^2(F_y)}$ poiché abbiamo usato la bilancia due volte. Otteniamo:
\begin{equation*}
u(\langle C_L \rangle)=0.0574 \qquad \Longrightarrow \qquad \langle C_L \rangle = -0.0280 \pm 0.0574\,.
\end{equation*}

Notiamo due cose:
\begin{enumerate}
    \item Il coefficiente di drag presenta un’incertezza maggiore, riflettendo la maggiore sensibilità di \( \langle C_D \rangle \) rispetto alla forza e alla portata \( Q \).
    \item Il fatto che l'incertezza di $\langle C_L \rangle$ sia più grande del valore medio significa che il valore medio di $\langle C_L \rangle$ è statisticamente indistinguibile da zero. Questo conferma perfettamente la teoria per un cilindro simmetrico.
\end{enumerate}

\newpage

\section*{Quesito 5}

Usare la funzione MATLAB \texttt{fft\_of\_force} fornita per calcolare la FFT (Fast Fourier Transform) del segnale $F_L$ e trovare lo spettro di frequenza. Identificare la frequenza di \emph{vortex shedding}, $f_{\text{shedding}}$, come il picco principale dello spettro di frequenza, e calcolare il numero di Strouhal $\text{Sr} = f_{\text{shedding}} \cdot D / U_b$. Confrontarlo con il valore di riferimento della letteratura allo stesso numero di Reynolds per il caso \emph{unbounded}. Come soluzione di riferimento, può essere utilizzato il grafico fornito nell'articolo Fey et al. (1998) (\href{https://doi.org/10.1063/1.869675}{doi}), come riportato in Figura~\ref{fig:str}.

\begin{figure}[H]
\centering
\includegraphics[width = 0.6 \textwidth]{str}
\caption{Numero di Strouhal dall'articolo Fey et al. (1998).}
\label{fig:str}
\end{figure}

\NotionDivider

Utilizzando la serie temporale di $F_L$ e la relativa frequenza di campionamento $f_s$, possiamo identificare la frequenza dominante di \emph{vortex shedding} come il picco dello spettro di frequenza calcolato con la FFT (Fast Fourier Transform). Quest'ultimo è mostrato in Figura~\ref{fig:freL}, mentre la frequenza di shedding è
\begin{equation*}
f_{\text{shedding}} = 1.008~\text{Hz}\,.   
\end{equation*}

È interessante notare che il calcolo di tale frequenza potrebbe essere \emph{pericoloso} se fatto sulla forza orizzontale. Questo perché ogni vortice che si stacca, sia sopra che sotto, contribuisce ad aumentare il drag. Inoltre, per via della turbolenza, ci potrebbe essere più rumore sperimentale.

\begin{figure}[H]
\centering
\includegraphics[width = 0.5 \textwidth]{test5img/freL}
\caption{Spettro di frequenza di $F_L$.}
\label{fig:freL}
\end{figure}

Ora calcoliamo il numero di Strouhal:
\[
\text{Sr} = \frac{f_{\text{shedding}} \cdot D}{U_b}
\approx 0.1814\,.
\]

I risultati di Fey et al. mostrano che
\begin{equation*}
\text{Sr}_{\text{Fey}} = \text{Sr}^* + \frac{m}{\sqrt{\text{Re}}}=0.1932\,,\text{ con }\text{Sr}^*=0.1776 \text{ e }m=2.2023\,.
\end{equation*}

La differenza è dunque del 6\%. Tale errore, considerabile comunque accettabile, è dovuto a queste principali differenze:
\begin{itemize}
    \item \textbf{Dominio}: Nel nostro esperimento il flusso avviene in una condotta d'acqua, mentre i dati sperimentali sono stati ottenuti in galleria del vento. La condotta introduce l'effetto dei bordi del dominio, che non è presente in un wind tunnel.
    \item \textbf{Errori Numerici}: I calcoli numerici introducono una discreta incertezza nei risultati.
\end{itemize}

Confermiamo quindi la consistenza fisica dei nostri risultati.

\section*{Quesito 6}

Impostare lo stesso livello dell'acqua nei test statici e dinamici è importante per la misurazione accurata della forza dinamica. Fornire una stima degli errori nei coefficienti di forza $\langle C_D\rangle$ e $\langle C_L\rangle$ dovuti a un diverso livello dell'acqua nei test statici e dinamici (assumere $\Delta h=2$ mm, per esempio). Considerare che l'area della sezione d'acqua di ciascuna delle due \emph{endplates}, come osservata dall'alto, è un rettangolo di dimensioni $180~\text{mm} \times 2~\text{mm}$ (vedere Figura~\ref{fig:test5sketch}).

\NotionDivider

L'effetto di una variazione di 2~mm nel livello dell'acqua (\(\Delta h = 2 ~ \text{mm}\)) tra i test è stato analizzato. La principale conseguenza di questa discrepanza è una modifica della forza di galleggiamento che agisce sul cilindro, influenzando i coefficienti \(\langle C_D\rangle\) e \(\langle C_L\rangle\).

La procedura seguita è stata la seguente:
\begin{itemize}
    \item L'area della sezione d'acqua delle piastre è stata calcolata come \(A_{\text{endplate}} = 0.18 \cdot 0.002 ~ \text{m}^2\).

    Di conseguenza, la differenza in volume è
    \begin{equation*}
    \Delta V = 2 \Delta h \cdot A_{\text{endplate}}
    \end{equation*} 

    \item La nuova forza di galleggiamento dovuta alla variazione del livello dell'acqua è stata calcolata come:
    \[
    F_B^{\text{new}} = F_B + \rho \cdot g \cdot \Delta V\,,
    \]
    dove \(F_B\) è la forza di galleggiamento originale, \(\rho = 998 ~ \text{kg/m}^3\) è la densità del fluido e \(g = 9.81 ~ \text{m/s}^2\) è l'accelerazione gravitazionale.

    \item La velocità media aggiornata \(U_b^{\text{new}}\) è stata determinata utilizzando la geometria del canale e la portata:
    \[
    U_b^{\text{new}} = \frac{Q}{B \cdot (h + \Delta h)}\,.
    \]

    \item I coefficienti \(C_L\) e \(C_D\) sono stati calcolati come:
    \[
    \langle C_D^{\text{new}} \rangle = \frac{\langle F_D\rangle}{0.5 \cdot \rho \cdot (U_b^{\text{new}})^2 \cdot A}\,, \qquad
    \langle C_L^{\text{new}}\rangle = \frac{\langle F_L^{\text{new}}\rangle}{0.5 \cdot \rho \cdot (U_b^{\text{new}})^2 \cdot A},
    \]
    dove \(A = D \cdot L\) rappresenta l'area proiettata del cilindro, \(\langle F_D \rangle\) e \(\langle F_L^{\text{new}} \rangle\) sono  rispettivamente la forza di resistenza media di Reynolds e la forza di sollevamento media di Reynolds aggiornata.

    \item Gli errori relativi nei coefficienti sono stati calcolati come:
    \[
    \mathcal{E}_{C_D} = \frac{|\langle C_D\rangle - \langle C_D^{\text{new}}\rangle|}{\langle C_D\rangle} \cdot 100=0.89\%\,, \qquad
    \mathcal{E}_{C_L} = \frac{|\langle C_L\rangle - \langle C_L^{\text{new}}\rangle|}{|\langle C_L\rangle|} \cdot 100=81.77\%\,.
    \]
\end{itemize}

I risultati evidenziano che anche piccole variazioni nel livello dell'acqua possono influire significativamente sull'accuratezza del coefficiente di portanza \(C_L\). Questo effetto è dovuto alla dipendenza diretta della forza di portanza dalla correzione della galleggiabilità, che risulta sensibile alle variazioni nel livello dell'acqua. Tuttavia, è bene ricordare che il valore di $C_L$ è molto basso, di conseguenza anche piccole variazioni hanno un notevole impatto percentuale.


 Al contrario, il coefficiente di resistenza \(C_D\), come si aspetterebbe, è meno influenzato.





 \cleardoublepage



\end{document}